% concatenated from complete_arXiv_ver1.tex
\documentclass[12pt]{article}
\usepackage{amsmath,amssymb}
\usepackage{amscd}
\usepackage{comment}
\usepackage{color}
\usepackage{enumerate}
\usepackage{cases}
\usepackage{mathtools}
\usepackage[all]{xy}
\usepackage{amsthm}
\usepackage{pb-diagram}
\usepackage[top=25truemm,bottom=25truemm,left=15truemm,right=15truemm]{geometry}
%\numberwithin{equation}{section}
\theoremstyle{definition}
\newtheorem{defi}{Definition}%[section]
\newtheorem{theo}[defi]{Theorem}
\newtheorem{lemm}[defi]{Lemma}
\newtheorem{prop}[defi]{Proposition}
\newtheorem{cor}[defi]{Corollary}
\newtheorem{ex}[defi]{Example}
\newtheorem{rem}[defi]{Remark}
\newtheorem{que}[defi]{Question}
\newtheorem{cond}[defi]{Condition}
\renewcommand{\proofname}{\it Proof.}
\newcommand{\sslash}{{/\mkern-6mu/}}
%%%%%%%%%%%%%%%%%%%
\def\ad{{\rm ad}}
\def\Aut{{\rm Aut}}
\def\det{{\rm det}}
\def\div{{\rm div}}
\def\pr{{\rm pr}}
\def\Spin{{\rm Spin}}
\def\Sp{{\rm Sp}}
\def\SO{{\rm SO}}
\def\SL{{\rm SL}}
\def\SU{{\rm SU}}
\def\U{{\rm U}}
\def\GL{{\rm GL}}
\def\O{{\rm O}}
\def\Ad{{\rm Ad}}
\def\End{{\rm End}}
\def\Exp{{\rm Exp}}
\def\skew{{\rm skew}}
\def\Exp{{\rm Exp}}
\def\Im{{\rm Im}}
\def\Gr{{\rm Gr}}
\def\Sym{{\rm Sym}}
\def\Hom{{\rm Hom}}
\def\coker{{\rm coker}}
\def\ind{{\rm ind}}
\def\det{{\rm det}}
\def\id{{\rm id}}
\def\Vol{{\rm Vol}}
\def\Re{{\rm Re}}
\def\Tr{{\rm Tr}}
\def\refe{\ast}
\def\ini{{\rm in}}
\def\vol{{\rm vol}}
\def\can{{\rm can}}
\def\max{{\rm max}}
\def\min{{\rm min}}

\def\H{{\mathbb H}}
\def\R{{\mathbb R}}
\def\Z{{\mathbb Z}}
\def\C{{\mathbb C}}
\def\N{{\mathbb N}}
\def\Q{{\mathbb Q}}
\def\K{{\mathbb K}}
\def\CP{{\mathbb C}{\mathbb P}}
\def\D{{\mathbb D}}


\def\sp{\mathfrak{sp}}
\def\so{\mathfrak{so}}
\def\spin{\mathfrak{spin}}
\def\g{\mathfrak{g}}
\def\t{\mathfrak{t}}
\def\h{\mathfrak{h}}
\def\s{\mathfrak{s}}
\def\k{\mathfrak{k}}
\def\u{\mathfrak{u}}

\def\A{{\cal A}}
\def\M{{\cal M}}
\def\F{{\cal F}}


\def\inum{{\sqrt{-1}}}

\providecommand{\keywords}[1]
{
\small 
\textbf{{Keywords---}} #1
}
\providecommand{\MSC}[1]
{
\small 
\textbf{{MSC---}} #1
}



\begin{document}

\title {Complete harmonic metrics and subharmonic functions \\
on the unit disc}
\author {Natsuo Miyatake}
\date{}
\maketitle
\begin{abstract} 
Let $X$ be a Riemann surface, $K_X \rightarrow X$ the canonical bundle, and $T_X\coloneqq K_X^{-1}\rightarrow X$ the dual bundle of the canonical bundle. For each integer $r \geq 2$, each $q \in H^0(K_X^r)$, and each choice of the square root $K_X^{1/2}$ of the canonical bundle, we canonically obtain a Higgs bundle $(\K_r,\Phi(q))\rightarrow X$, which is called a cyclic Higgs bundle. A diagonal harmonic metric $h = (h_1, \dots, h_r)$ on a cyclic Higgs bundle yields $r-1$-Hermitian metrics $H_1, \dots, H_{r-1}$ on $T_X\rightarrow X$, defined as $H_j \coloneqq h_j^{-1} \otimes h_{j+1}$ for each $j=1,\dots, r-1$, while $h_1$, $h_r$, and $q$ yield a degenerate Hermitian metric $H_r$ on $T_X \rightarrow X$. A diagonal harmonic metric is said to be complete if the K\"ahler metrics induced by $H_1,\dots, H_{r-1}$ are all complete. Li-Mochizuki established a theorem stating that on any Riemann surface $X$ and any $q$ that is non-zero unless $X$ is hyperbolic, there exists a unique complete harmonic metric $h$ on $(\K_r,\Phi(q))\rightarrow X$ with a fixed determinant. The holomorphic section $q$ induces a subharmonic weight function $\phi_q=\frac{1}{r}\log|q|^2$ on $K_X\rightarrow X$, and a diagonal harmonic metric depends solely on this weight function $\phi_q$. In this paper, we extend the uniqueness part of the theorem of Li-Mochizuki to any subharmonic weight function $\varphi$ whose exponential is $C^2$ outside a compact subset $K \subseteq X$. 
We also show that on the unit disc $\D \coloneqq \{z \in \C \mid |z| < 1\}$, a complete Hermitian metric associated with $\varphi$ always exists. 
Furthermore, on the unit disc, when $\varphi$ can be monotonically approximated by a family of weight functions $(\varphi_\epsilon)_{0 < \epsilon < 1}$, where each $\varphi_\epsilon$ is smooth and defined on a disc $\D_\epsilon \coloneqq \{z \in \C \mid |z| < 1 - \epsilon\}$, we show that the corresponding family of complete metrics $(h_\epsilon)_{0 < \epsilon < 1}$ converges monotonically to a complete metric $h$ associated with $\varphi$ as $\epsilon\searrow 0$.

\end{abstract}


\MSC{30C15, 31A05, 53C07}

\keywords{Cyclic Higgs bundles, Harmonic metrics, Complete solutions, Subharmonic functions}

\section{Introduction and main theorem}\label{1}
Li-Mochizuki \cite{LM1} established a theorem stating that for any $r$-differential $q \in H^0(K_X^r)$ with $r\geq 2$, which is non-zero unless $X$ is hyperbolic, there uniquely exists a complete harmonic metric with fixed determinant on the cyclic Higgs bundle associated with $q$.
The purpose of this paper is to extend the uniqueness part of this theorem to Hermitian metrics associated with more general subharmonic weight functions satisfying a certain regularity condition, and to establish both existence and approximation theorems for complete Hermitian metrics associated with subharmonic weight functions on the unit disc. The objective of this section is to state the concrete assertion of our main theorem. Before presenting the main theorem, we explain the background of cyclic Higgs bundles, harmonic metrics, and their extensions using subharmonic weight functions.

Let $X$ be a connected Riemann surface and $K_X\rightarrow X$ the canonical bundle. We choose a square root $K_X^{1/2}$ of the canonical bundle $K_X\rightarrow X$. Let $\K_r$ be a holomorphic vector bundle of rank $r$ defined as follows:
\begin{align*}
\K_r\coloneqq \bigoplus_{j=1}^r K_X^{\frac{r-(2j-1)}{2}} = K_X^{\frac{r-1}{2}} \oplus K_X^{\frac{r-3}{2}} \oplus \cdots \oplus K_X^{-\frac{r-3}{2}} \oplus K_X^{-\frac{r-1}{2}}.
\end{align*}
For each $q\in H^0(K_X^r)$, we associate a holomorphic section $\Phi(q)\in H^0(\End \K_r\otimes K_X)$, which is called a Higgs field, as follows:
\begin{align*}
\Phi(q)\coloneqq 
\left(
\begin{array}{cccc}
0 & && q \\ 
1 & \ddots && \\ 
& \ddots & \ddots & \\ 
& & 1 & 0
\end{array}
\right).
\end{align*}
The pair $(\K_r,\Phi(q))\rightarrow X$ is called a cyclic Higgs bundle (cf. \cite{Bar1}). Since the cyclic Higgs bundle $(\K_r,\Phi(q))\rightarrow X$ is a Higgs bundle, it is of course related to the general theory of Higgs bundles, but in this paper we would like to emphasize the view that a harmonic metric on it provides an extended concept of non-negative constant Gaussian curvature K\"ahler metrics, as the author has stated in other papers \cite{Miy4, Miy5}. The connection between harmonic metrics on cyclic Higgs bundles and the non-negative constant curvature K\"ahler metrics is as follows: We consider the case where $r=2$ for simplicity. This is the most famous example of Higgs bundles, which was introduced by Hitchin \cite[Section 11]{Hit1}. Suppose that $X$ is a compact Riemann surface. Firstly, if $X=\CP^1$, then $(\K_2,\Phi(q)=\Phi(0))\rightarrow X$ is unstable, and thus there are no solutions to the Hitchin equation. Secondly, if $X$ is an elliptic curve, then there exists a solution to the Hitchin equation for $(\K_2,\Phi(q))\rightarrow X$ if and only if $q\neq 0$. If $q$ is a non-zero section, then we can easily find a solution to the Hitchin equation; $h=(h_1, h_1^{-1})=(h_q^{-1/4}, h_q^{1/4})$ solves the Hitchin equation, where $h_q$ is a Hermitian metric on $K_X^{-2}\rightarrow X$ defined as follows:
\begin{align}
(h_q)_x(u,v)\coloneqq \langle q_x,u\rangle\overline{\langle q_x,v\rangle} \ \text{for $x\in X$, $u,v\in (K_X^{-2})_x$}. \label{uv}
\end{align}
Finally, if $X$ is a compact hyperbolic Riemann surface, then $(\K_2,\Phi(q))\rightarrow X$ is always stable, and thus there uniquely exists a harmonic metric $h$ on $(\K_2, \Phi(q))\rightarrow X$ satisfying $\det(h)=1$. Moreover, from the symmetry of the Higgs field and the uniqueness of the solution to the Hitchin equation, the metric $h$ splits as $h=(h_1,h_1^{-1})$. As described above, when $X$ is an elliptic curve and $q\neq0$, or when $X$ is a hyperbolic Riemann surface, there exists a diagonal harmonic metric on $(\K_2,\Phi(q))\rightarrow X$. A diagonal harmonic metric naturally induces a Hermitian metric on $K_X^{-1}\rightarrow X$ defined as $H_1\coloneqq h_1^{-2}$. Furthermore, by combining it with the Hermitian metric $h_q$ induced by $q$, we obtain a Hermitian metric $H_2$ on $K_X^{-1}\rightarrow X$ which is degenerated at the zeros of $q$ as $H_2\coloneqq h_q\otimes H_1^{-1}=h_q\otimes h_1^2$. If $X$ is an elliptic curve, then we have $H_1=H_2=h_q^{1/2}$ and it induces a flat Gaussian curvature K\"ahler metric. When $X$ is a compact hyperbolic Riemann surface, the harmonic metric is distorted by the zero point of $q$ and deviates from a flat metric, and in the ``fully distorted'' limit, i.e., when $q=0$, the K\"ahler metric induced by $H_1$ becomes a metric of negative constant curvature. The above is an explanation for the case where $r=2$, but the same thing can be observed for harmonic metrics on higher order cyclic Higgs bundles (see Section \ref{2}).

When the harmonic metric on a cyclic Higgs bundle is regarded as an extension of the non-negative constant curvature K\"ahler metric, the additional structure $q$ that causes the deviation from the hyperbolic metric does not need to be a holomorphic section. In the papers \cite{Miy3, Miy4, Miy5}, the author proposed a new research direction to extend the singular Hermitian metric $h_q^{-1/r}$ induced by $q$ to more general semipositive singular Hermitian metric $e^{-\varphi}h_\refe$ on $K_X\rightarrow X$. This means that we consider harmonic metrics on cyclic ramified coverings of $X$ and the limit of the covering order going to infinity (see Section \ref{2.2}). The aim is to deepen the integration of classical potential theory and complex analysis, especially mathematics related to the zeros of holomorphic functions and holomorphic sections (cf. \cite{BCHM1, Ran1, ST1, SZ1}), with harmonic metrics. For more specific research results at this stage, see \cite{Miy4, Miy5}.

A diagonal harmonic metric $h = (h_1, \dots, h_r)$ on a cyclic Higgs bundle yields $r-1$-Hermitian metrics $H_1, \dots, H_{r-1}$ on $T_X\rightarrow X$, defined as $H_j \coloneqq h_j^{-1} \otimes h_{j+1}$ for each $j=1,\dots, r-1$, while $h_1$, $h_r$, and $q$ yield a degenerate Hermitian metric $H_r\coloneqq h_r^{-1}\otimes h_1\otimes h_q$ on $T_X \rightarrow X$. A diagonal harmonic metric is said to be {\it complete} if the K\"ahler metrics induced by $H_1,\dots, H_{r-1}$ are all complete. Li-Mochizuki \cite{LM1} established a theorem stating that on any open Riemann surface $X$ and any $q$ that is non-zero unless $X$ is hyperbolic, there exists a unique complete harmonic metric $h$ with a fixed determinant. The notion of complete harmonic metrics is naturally extended to general semipositive singular Hermitian metrics (see Sections \ref{rdc} and \ref{2.2}). The purpose of this paper is to extend the uniqueness part of this theorem to Hermitian metrics associated with more general semipositive singular metric $e^{-\varphi}h_\refe$ on $K_X\rightarrow X$ satisfying a certain regularity condition, and to establish both existence and approximation theorems for complete Hermitian metrics associated with $e^{-\varphi}h_\refe$ on the unit disc. In the following, when we say ``$h$ is a complete solution associated with $\varphi$", we mean that $h$ is a solution to the extended Hitchin equation (\ref{phi}) in Section \ref{2.2} associated with $\varphi$ and it is complete in the sense described in Section \ref{rdc}. Our main theorems are as follows:

\begin{theo}\label{main theorem 1} {\it Suppose that $\varphi$ is not identically $-\infty$ unless $X$ is hyperbolic. Suppose also that $\varphi$ satisfies the following assumption (=Condition \ref{condition} in Section \ref{4}):
\begin{enumerate}[$(\ast)$]
\item There exists a compact subset $K\subseteq X$ such that on $X\backslash K$, $e^{\varphi}$ is of class $C^2$.
\end{enumerate}
Then for any two complete solutions $h=(h_1,\dots, h_r)$ and $h^\prime=(h_1^\prime,\dots, h^\prime)$ associated with $\varphi$ satisfying $\det(h)=\det(h^\prime)=1$, we have $h=h^\prime$. 
}
\end{theo}

\begin{theo}\label{main theorem 2} {\it On the unit disc $\D\coloneqq \{z\in\C\mid |z|<1\}$, for any $\varphi$, there exists a real complete solution $h=(h_1,\dots, h_r)$ associated with $\varphi$. For each $j=1,\dots, [r/2]$, $h_j$ is of class $C^{2j-1,\alpha}$ for any $\alpha\in(0,1)$. }
\end{theo}
\begin{theo}\label{main theorem 3}{\it
Let $(e^{-\varphi_\epsilon}h_\refe)_{\epsilon>0}$ be a family of smooth semipositive Hermitian metrics on the canonical bundle, each of which is defined on a disc $\D_\epsilon \coloneqq \{z \in \C \mid |z| < 1 - \epsilon\}$, that satisfies the following property:
\begin{itemize}
\item For each $\epsilon>\epsilon^\prime>0$, $\varphi_{\epsilon^\prime}\leq \varphi_\epsilon$ and $(\varphi_\epsilon)_{0<\epsilon<1}$ converges to a function $\varphi$ on $\D$ that is locally a sum of a smooth function and a subharmonic function as $\epsilon\searrow 0$.
\end{itemize}
Then the corresponding family of smooth complete solutions $(h_\epsilon=(h_{1,\epsilon},\dots, h_{r,\epsilon}))_{\epsilon>0}$ monotonically converges to a complete solution $h=(h_1,\dots, h_r)$ associated with $\varphi$ as $\epsilon\searrow 0$.}
\end{theo}

The assumption $(\ast)$ in Theorem \ref{main theorem 1} is for using the maximum principle on open Riemann surfaces (cf. \cite[Section 3]{LM1}). The author expects that this assumption is unnecessary, but has not yet been able to prove the uniqueness without it. If we can drop the assumption $(\ast)$ in Theorem \ref{main theorem 1}, then, it follows from Theorem \ref{main theorem 2} that for every $\varphi$ on a hyperbolic surface, by lifting to the universal covering space and applying the uniqueness argument on the unit disc (cf. \cite[Proposition 5.7]{LM1}), there exists a complete solution associated with $\varphi$.

The approximation $(\varphi_\epsilon)_{0<\epsilon<1}$ of $\varphi$ in Theorem \ref{main theorem 3} always exists by the standard theory of mollification of subharmonic functions (cf. \cite{Ran1}). For the proof of Theorem \ref{main theorem 3}, we use the technique established in \cite[Proposition 4.6]{LM1}. This approximation is used in \cite{Miy4, Miy5} to prove theorems about functions called entropy and free energy, which are defined using $h=(h_1,\dots, h_r)$ and $e^{-\varphi}h_\refe$. The analogy of Theorems \ref{main theorem 2} and \ref{main theorem 3} on the complex plane will be addressed in a subsequent paper.

The structure of this paper is as follows: In Section \ref{2}, we will explain some backgrounds about harmonic metrics on cyclic Higgs bundles. In Section \ref{3}, we will review the Omori-Yau maximum principle and the Cheng-Yau maximum principle which will be used in the proof of Theorem \ref{main theorem 1}. In Section \ref{4}, we will give the proofs of Theorems \ref{main theorem 1}, \ref{main theorem 2}, and \ref{main theorem 3}.

\noindent
{\bf Acknowledgements.} I would like to thank Toshiaki Yachimura for helpful conversations. This work was supported by the Grant-in-Aid for Early Career Scientists (Grant Number: 24K16912).

\section{Backgrounds}\label{2}
\subsection{Real, diagonal, and complete Hermitian metrics}\label{rdc}
Let $X$ be a connected, possibly non-compact Riemann surface and $K_X\rightarrow X$ the canonical bundle. We choose a square root $K_X^{1/2}$ of the canonical bundle $K_X\rightarrow X$. Let $\K_r$ be a holomorphic vector bundle of rank $r$ defined as follows:
\begin{align*}
\K_r\coloneqq \bigoplus_{j=1}^r K_X^{\frac{r-(2j-1)}{2}} = K_X^{\frac{r-1}{2}} \oplus K_X^{\frac{r-3}{2}} \oplus \cdots \oplus K_X^{-\frac{r-3}{2}} \oplus K_X^{-\frac{r-1}{2}}.
\end{align*}
A Hermitian metric $h$ on $\K_r\rightarrow X$ is said to be {\it real} (cf. \cite{Hit2}) if the following bundle isomorphism $S$ is isometric with respect to $h$:
\begin{align*}
S\coloneqq \left(
\begin{array}{ccc}
& & 1 \\ 
& \reflectbox{$\ddots$} & \\ 
1 & &
\end{array}
\right): \K_r\rightarrow \K_r^{\vee},
\end{align*}
where $\K_r^{\vee}$ denotes the dual bundle of $\K_r$. Also, $h$ is called {\it diagonal} (cf. \cite{Bar1, Col1, DL2}) if it splits as $h=(h_1,\dots, h_r)$. Let $T_X\rightarrow X$ be the dual bundle of $K_X\rightarrow X$. By taking the difference between the adjacent components in a diagonal Hermitian metric $h=(h_1,\dots, h_r)$, we obtain $r-1$-Hermitian metrics $H_1,\dots, H_{r-1}$ on $T_X\rightarrow X$:
\begin{align*}
H_j\coloneqq h_j^{-1}\otimes h_{j+1}.
\end{align*}
A diagonal Hermitian metric $h$ is real if and only if $h_j=h_{r-j+1}$ for all $j=1,\dots, r$. This is also equivalent to the metrics $H_1,\dots, H_{r-1}$ satisfying $H_j=H_{r-j}$ for all $j=1,\dots, r-1$. We say that a diagonal Hermitian metric $h$ is {\it complete} (cf. \cite{LM1}) if the K\"ahler metrics induced by $H_1,\dots, H_{r-1}$ are all complete. 
\subsection{Harmonic metrics on cyclic Higgs bundles}\label{2.1}
Let $q$ be a holomorphic section of $K_X^r\rightarrow X$. Then we associate a holomorphic section $\Phi(q)\in H^0(\End \K_r\otimes K_X)$, which is called a Higgs field, as follows:
\begin{align*}
\Phi(q)\coloneqq 
\left(
\begin{array}{cccc}
0 & && q \\ 
1 & \ddots && \\ 
& \ddots & \ddots & \\ 
& & 1 & 0
\end{array}
\right).
\end{align*}
The pair $(\K_r,\Phi(q))$ is called a cyclic Higgs bundle (cf. \cite{Bar1, Hit1, Hit2}). The above cyclic Higgs bundle is an example of cyclotomic Higgs bundles introduced in \cite{Sim3}. For each Hermitian metric $h$ on $\K_r$, we associate a connection $D_q(h)$ defined as follows:
\begin{align*}
D_q(h)\coloneqq \nabla^h + \Phi(q) + \Phi(q)^{\ast h}.
\end{align*}
A Hermitian metric on $(\K_r,\Phi(q))\rightarrow X$ is called a {\it harmonic metric} if the connection $D_q(h)$ is a flat connection. A harmonic metric $h$ yields a harmonic map $\hat{h}: \widetilde{X}\rightarrow \SL(r,\C)/\SU(r)$, where $\widetilde{X}$ is the universal covering space. If $h$ is real, then $\hat{h}$ maps into $\SL(r,\R)/\SO(r)$. Suppose that $X$ is a compact Riemann surface of genus at least $2$. Then the Higgs bundle $(\K_r,\Phi(q))$ is stable for any $q\in H^0(K_X^r)$, and thus there uniquely exists a harmonic metric $h$ on $(\K_r,\Phi(q))$ such that $\det(h)=1$. From the uniqueness of the metric and the symmetry of the Higgs field (cf. \cite{Bar1, Hit2}), the harmonic metric $h$ is diagonal and real. A non-diagonal harmonic metric (resp. a non-real harmonic metric) can be constructed by considering the Dirichlet problem (cf. \cite{Don1, LM2}) for the Hitchin equation with a non-diagonal (resp. a non-real) Hermitian metric on the boundary. 
Let $h=(h_1,\dots, h_r)$ be a diagonal harmonic metric on a cyclic Higgs bundle $(\K_r,\Phi(q))\rightarrow X$. As mentioned in the previous subsection we obtain $r-1$-Hermitian metrics $H_1,\dots, H_{r-1}$ on $T_X\rightarrow X$ defined as
\begin{align*}
H_j\coloneqq h_j^{-1}\otimes h_{j+1} \ \text{for each $j=1,\dots,r -1$}.
\end{align*}
We also obtain a Hermitian metric $H_r$ on $T_X\rightarrow X$ which is degenerate at the zeros of $q$ as follows:
\begin{align*}
H_r\coloneqq h_r^{-1}\otimes h_1 \otimes h_q,
\end{align*}
where $h_q$ is a Hermitian metric on $K_X^{-r}\rightarrow X$ that is degenerate at the zeros of $q$ defined as follows:
\begin{align}
(h_q)_x(u,v)\coloneqq \langle q_x,u\rangle\overline{\langle q_x,v\rangle} \ \text{for $x\in X$, $u,v\in (K_X^{-r})_x$}. \label{q}
\end{align}
The connection $D_q(h)$ is flat if and only if the Hermitian metric $h$ satisfies the following elliptic equation, called the Hitchin equation:
\begin{align*}
F_h + [\Phi(q)\wedge\Phi(q)^{\ast h}] = 0.
\end{align*}
By using $H_1,\dots, H_r$, we can describe the Hitchin equation as follows: 
\begin{align}
\inum F_{h_j} + \vol(H_{j-1}) - \vol(H_j) = 0 \ \text{for $j=1,\dots, r-1$}, \label{F}
\end{align}
where for each $j=1,\dots, r$, $\vol(H_j)$ is the volume form of the metric $H_j$ which is locally described as $\vol(H_j)=\inum H_j(\frac{\partial}{\partial z}, \frac{\partial}{\partial z})\ dz\wedge d\bar{z}$, and $H_0$ is $H_r$. Equation (\ref{F}) is a kind of Toda equation (cf. \cite{AF1, Bar1, GH1, GL1, Moc0, Moc1}). By making a variable change from $h=(h_1,\dots, h_r)$ to $H_1,\dots, H_{r-1}$, the Hitchin equation becomes the following:
\begin{align}
\inum F_{H_j} + 2\vol(H_j) - \vol(H_{j-1}) - \vol(H_{j+1}) = 0 \ \text{for $j=1,\dots, r-1$.} \label{cyc}
\end{align}
In the case where $r=2$ and $q=0$, the Hitchin equation for the cyclic Higgs bundle is equivalent to the following equation:
\begin{align}
\inum F_H + 2\vol(H) = 0, \label{H}
\end{align}
where the solution $H$ to equation (\ref{H}) is a Hermitian metric $H$ on $T_X\rightarrow X$. The K\"ahler metric induced by $H$ is a constant negative Gaussian curvature K\"ahler metric (cf. \cite{Hit1}). 
We remark on the following two points:
\begin{itemize}
\item If $q$ has no zeros, then $H_1=\cdots=H_{r-1}=h_q^{1/r}$ is a solution to PDE (\ref{cyc}).
\item Suppose that $q=0$ and that there is a solution $H$ to equation (\ref{H}). We define positive constants $\lambda_1,\dots, \lambda_{r-1}$ as follows: 
\begin{align}
\lambda_j\coloneqq 2\sum_{k=1}^{r-1}(\Lambda_{r-1}^{-1})_{jk} = j(r-j), \label{lambda}
\end{align}
where $\Lambda_{r-1}^{-1}$ is the inverse matrix of the Cartan matrix of type $A_{r-1}$. Then $H_j=\lambda_j H\ (j=1,\dots, r-1)$ is a solution to PDE (\ref{cyc}).
\end{itemize}

Li-Mochizuki \cite{LM1} established the following result:
\begin{theo}[\cite{LM1}]{\it Suppose that $q$ is a non-zero holomorphic section unless $X$ is hyperbolic. Then there always exists a unique complete harmonic metric $h$ on $(\K_r,\Phi(q))\rightarrow X$ satisfying $\det(h)=1$. Moreover, the complete harmonic metric $h$ is real.
}
\end{theo}
In addition to \cite{LM1}, we refer the reader to \cite{DL3, DW1, GL1, LT1, LTW1, Li1, LM2, Moc0, Moc1, Nie1, Wan1, WA1} for the works related to complete harmonic metrics. 

\subsection{More general subharmonic weight functions}\label{2.2}
For each $q\in H^0(K_X^r)$, the metric $h_q$ induces a singular metric (cf. \cite{GZ1}) $h_q^{-1/r}$ that diverges at the zeros of $q$. We can define a curvature for $h_q^{-1/r}$ in the sense of currents, which is semipositive and has support at the zero points of $q$. We fix a smooth metric $h_\refe$ and denote by $e^{-\phi_q}h_\refe$ the metric $h_q^{-1/r}$, where the weight function $\phi_q$ is defined as $\phi_q\coloneqq \frac{1}{r}\log|q|_{h_\refe}^2$. If we choose a local flat reference metric, then the weight function defines a local subharmonic function. Also, although we will not go into details here, there is a way to identify the weight function with a global psh function on the total space of the dual line bundle excluding zero points, see \cite[Section 2.1]{BB1}. In the sense described above, we call the function $\phi_q$ a {\it subharmonic weight function}, although this is a bit of an abuse of terminology. We consider a more general singular metric $e^{-\varphi}h_\refe$ with semipositive curvature, where the weight function $\varphi$ is a function that is locally a sum of a subharmonic function and a smooth function. For each $q_N\in H^0(K_X^N)$, let $\phi_{q_N}$ be defined as $\frac{1}{N}\log|q_N|_{h_\refe}^2$. Then $e^{-\phi_{q_N}}h_\refe$ is a singular metric with semipositive curvature. Moreover, any semipositive singular metric can be approximated, at least in the $L^1_{loc}$ sense, by a sequence $(e^{-\phi_{q_N}}h_\refe)_{N\in\N}$, where $q_N\in H^0(K_X^N)$ (cf. \cite{GZ1}). In \cite{Miy2, Miy3}, for each semipositive singular metric $e^{-\varphi}h_\refe$ and each $r\geq 2$, the following equation for a diagonal metric $h=(h_1,\dots, h_r)$ on $\K_r\rightarrow X$ was introduced:
\begin{align}
\inum F_{h_j} + \vol(H_{j-1}) - \vol(H_j) = 0 \ \text{for $j=1,\dots, r-1$}, \label{phi}
\end{align}
where $H_1,\dots, H_{r-1}$ are defined as $H_j\coloneqq h_j^{-1}\otimes h_{j+1}$ for each $j=1,\dots, r-1$, and $H_r=H_0$ is defined as follows:
\begin{align*}
H_r = H_0 &\coloneqq h_r^{-1}\otimes h_1 \otimes (e^{-\varphi}h_\refe)^{-r} \\ 
&= H_1^{-1}\otimes \cdots \otimes H_{r-1}^{-1} \otimes (e^{-\varphi}h_\refe)^{-r}.
\end{align*}
For the case where $\varphi=\phi_q=\frac{1}{r}\log|q|_{h_\refe}^2$, equation (\ref{phi}) becomes the Hitchin equation for the cyclic Higgs bundle $(\K_r,\Phi(q))$. For the case where $\varphi=\phi_{q_N}=\frac{1}{N}\log|q_N|_{h_\refe}^2$, equation (\ref{phi}) gives a harmonic metric on a ramified covering space of $X$ (see \cite[Section 2]{Miy2}). As mentioned above, more general $\varphi$ can be approximated by a sequence $(\phi_{q_N})_{N\in\N}$ at least in the $L^1_{loc}$-sense (cf. \cite{GZ1}), therefore equation (\ref{phi}) can be considered to be an equation obtained as the limit when the covering degree increases infinitely or the number of zeros increases infinitely. By making a variable change from $h=(h_1,\dots, h_r)$ to $H_1,\dots, H_{r-1}$, the equation becomes the following:
\begin{align}
\inum F_{H_j} + 2\vol(H_j) - \vol(H_{j-1}) - \vol(H_{j+1}) = 0 \ \text{for $j=1,\dots, r-1$.}\label{cyc2}
\end{align}
We remark on the following two points:
\begin{itemize}
\item If $e^{-\varphi}h_\refe$ is flat, then $H_1=\cdots=H_{r-1}=(e^{-\varphi}h_\refe)^{-1}$ is a solution to PDE (\ref{cyc2}).
\item Suppose that $\varphi=-\infty$. Then equation (\ref{cyc2}) is the same as equation (\ref{cyc}) for the case where $q=0$, and thus for a solution $H$ to equation (\ref{H}), $H_j=\lambda_j H\ (j=1,\dots, r-1)$ is a solution to PDE (\ref{cyc2}), where $\lambda_1,\dots, \lambda_{r-1}$ are constants defined in (\ref{lambda}).
\end{itemize}






\subsection{Entropy and free energy}\label{S and F}
This subsection reviews on the definitions of the functions called entropy and free energy given in \cite{Miy4, Miy5}. Let $h=(h_1,\dots, h_r)$ be a solution to equation (\ref{phi}) and let $H_1,\dots, H_r$ be the Hermitian metrics constructed from $h$ and $e^{-\varphi}h_\refe$ as in Section \ref{2.2}. Then we define functions, which we call entropy and free energy, as follows:
\begin{defi}\label{entropy2} 
For each $j=0, 1,\dots, r-1$ and non-zero real number $\beta$, let $p_j(r,\beta, \varphi):X\rightarrow [0,1]$ be a nonnegative function defined as follows:
\begin{align*}
p_j(r,\beta, \varphi)\coloneqq \frac{\vol(H_j)^\beta}{\sum_{j=0}^{r-1}\vol(H_j)^\beta},
\end{align*}
where $\vol(H_j)^\beta/\sum_{j=0}^{r-1}\vol(H_j)^\beta$ is understood to be $(\vol(H_j)/\vol(H_\refe))^\beta/\sum_{j=0}^{r-1}(\vol(H_j)/\vol(H_\refe))^\beta$ for some reference metric $H_\refe$, which does not affect $p_j(r,\beta,\varphi)$. We call the following function {\it entropy}:
\begin{align*}
S(r,\beta, \varphi)\coloneqq -\sum_{j=0}^{r-1} p_j(r, \beta, \varphi)\log p_j(r,\beta, \varphi).
\end{align*}
\end{defi}
\begin{defi}\label{F} Let $\beta$ be a non-zero real number. We call the following function {\it free energy}:
\begin{align*}
F(r,\beta,\varphi, H_\refe)\coloneqq -\frac{1}{\beta}\log(\sum_{j=0}^{r-1}(\vol(H_j)/\vol(H_\refe))^\beta),
\end{align*}
where $H_\refe$ is a reference metric. 
\end{defi}
Note that the difference between free energy functions with the same reference metric does not depend on the choice of the reference metric.

\section{The Omori-Yau maximum principle and the Cheng-Yau maximum principle}\label{3}
This section reviews the Omori-Yau maximum principle \cite{Omo1, Yau1} and the Cheng-Yau maximum principle \cite{CY1}, along with their variants for functions on manifolds with boundary \cite{LM1}. The following is the Omori-Yau maximum principle \cite{Omo1, Yau1} (see also \cite[Lemma 3.1]{LM1}):
\begin{theo}[\cite{Omo1, Yau1}]{\it Let $(M, g_M)$ be a complete Riemannian manifold with Ricci curvature bounded from below. Then for each real-valued $C^2$-function $u$ on $M$ that is bounded above, there exist an integer $m_0 \in \Z_{\geq 1}$ and a sequence of points $(x_m)_{m \geq m_0}$ in $M$ such that for each $m \geq m_0$,
\begin{align*}
&u(x_m) \geq \sup_M u - \frac{1}{m}, \\
&|du|_{g_M}(x_m) \leq \frac{1}{m}, \\
&(\Delta_{g_M} u)(x_m) \leq \frac{1}{m},
\end{align*}
where $|du|_{g_M}$ denotes the norm of the exterior derivative of $u$, and where $\Delta_{g_M} := -d^{\ast}d$ denotes the negative Laplacian with respect to $g_M$.
}
\end{theo}

The following is the Cheng-Yau maximum principle \cite{CY1} (see also \cite[Lemma 3.3]{LM1}):

\begin{theo}[\cite{CY1}]\label{C-Ym}{\it Let $(M, g_M)$ be a complete Riemannian manifold with Ricci curvature bounded from below. Let $u$ be a real-valued $C^2$-function on $M$ satisfying $\Delta_{g_M} u \geq f(u)$, where $f:\R\rightarrow \R$ is a function. Suppose that there exists a continuous positive function $g:[a,\infty)\rightarrow \R_{>0}$ such that
\begin{enumerate}[(i)]
\item $g$ is non-decreasing;
\item $\liminf_{t\to\infty}\frac{f(t)}{g(t)}>0$;
\item $\int_a^\infty(\int_b^tg(\tau)d\tau)^{-1/2}dt<\infty$ for some $b\geq a$.
\end{enumerate}
Then the function $u$ is bounded above. Moreover, if $f$ is lower semicontinuous, then $f(\sup_M u)\leq 0$.
}
\end{theo}

In the proof of Theorem \ref{main theorem 1}, we will use the following variants of the Omori-Yau maximum principle and the Cheng-Yau maximum principle for functions on manifolds with boundary, which were established in \cite[Lemmas 3.2 and 3.4]{LM1}.

\begin{theo}[\cite{LM1}]\label{OYMB}
{\it Let $(M, g_M)$ be a complete Riemannian manifold with smooth compact boundary $\partial M$. Suppose that the Ricci curvature of $g_M$ is bounded from below. Then for each real-valued $C^2$-function $u$ on $M$ that is bounded above, either of the following holds:
\begin{itemize}
\item $\sup_{M} u = \max_{\partial M} u$;
\item There exist an integer $m_0 \in \Z_{\geq 1}$ and a sequence of points $(x_m)_{m \geq m_0}$ in $M$ such that for each $m \geq m_0$,
\begin{align*}
&u(x_m) \geq \sup_M u - \frac{1}{m}, \\
&|du|_{g_M}(x_m) \leq \frac{1}{m}, \\
&(\Delta_{g_M} u)(x_m) \leq \frac{1}{m}.
\end{align*}
\end{itemize}
}
\end{theo}

\begin{theo}[\cite{LM1}]\label{CYMB}
{\it Let $(M, g_M)$ be a complete Riemannian manifold with smooth compact boundary $\partial M$. Suppose that the Ricci curvature of $g_M$ is bounded from below. Let $u$ be a real-valued $C^2$-function on $M$ satisfying $\Delta_{g_M} u \geq f(u)$, where $f:\R\rightarrow \R$ is a function. Suppose that there exists a continuous positive function $g:[a,\infty)\rightarrow \R_{>0}$ such that
\begin{enumerate}[(i)]
\item $g$ is non-decreasing;
\item $\liminf_{t\to\infty}\frac{f(t)}{g(t)} > 0$;
\item $\int_a^\infty \left( \int_b^t g(\tau)\, d\tau \right)^{-1/2} dt < \infty$ for some $b \geq a$.
\end{enumerate}
Then the function $u$ is bounded above. Moreover, if $f$ is lower semicontinuous, then one of the following holds:
\begin{itemize}
\item $\sup_M u = \max_{\partial M} u$;
\item $f(\sup_M u) \leq 0$.
\end{itemize}
}
\end{theo}

Note that in Theorems \ref{OYMB} and \ref{CYMB}, we have in mind non-compact manifolds with compact boundary, such as those obtained by removing a relatively compact open subset from a non-compact manifold.



\section{Proof}\label{4}
\subsection{Proof of Theorem \ref{main theorem 1}}
We establish Theorem \ref{main theorem 1}. As mentioned in Section \ref{1}, this paper does not address the uniqueness of a completely general weight function $\varphi$. Instead, we establish uniqueness under condition $(\ast)$ in Theorem \ref{main theorem 1}, which we restate below:

\begin{cond}\label{condition} There exists a compact subset $K\subseteq X$ such that on $X\backslash K$, $e^{\varphi}$ is of class $C^2$.
\end{cond} 
For the proof of the uniqueness of complete solutions to equation (\ref{phi}), we review the lemmas on linear algebra established in \cite[Section 3.2]{LM1}. Note that the symbols have been slightly altered from those used in \cite{LM1} for the convenience of our use. Let $h_1,\dots, h_r$ be positive real numbers satisfying $\Pi_{i=1}^r h_i=h_1\cdot h_2\cdots h_{r-1}\cdot h_r=1$. We set $H_j\coloneqq h_j^{-1}\cdot h_{j+1}$ for each $j=1,\dots, r-1$. Let $\phi$ be a real number or $-\infty$. We set a non-negative real number $H_r$ as $H_r\coloneqq H_1^{-1}\cdot H_2^{-1}\cdots H_{r-1}^{-1}\cdot e^{r\phi}=h_r^{-1}\cdot h_1\cdot e^{r\phi}$. We also set $H_0\coloneqq H_r$. The following Lemma \ref{lem1} and Lemma \ref{lem2} are special cases of well-known inequalities which are formulated in more general settings (cf. \cite{LM1, Sim2}).
\begin{lemm}[cf. \cite{LM1, Sim2}]\label{lem1}{\it There exists a positive constant $C$ which depends only on $r$ such that the following holds for any $h_1,\dots, h_r$ satisfying $\Pi_{i=1}^r h_i=h_1\cdot h_2\cdots h_{r-1}\cdot h_r=1$ and $\phi$:
\begin{align}
\sum_{j=1}^r(H_{j-1}-H_j)^2\geq C(\sum_{j=1}^rH_j-re^\phi)^2. \label{ephi1}
\end{align}
}
\end{lemm}
\begin{proof} We give a proof directly without using the general theory. Since the variable transformation from $h_1,\dots, h_{r-1}$ to $H_j=h_j^{-1}\cdot h_{j+1}$ is inversible, it is sufficient to state that inequality (\ref{ephi1}) holds for any positive real numbers $H_1,\dots, H_{r-1}$. We first suppose that $\phi$ is not $-\infty$. Note that it is enough to consider the case where $\phi=0$, in other words, $e^\phi=1$. This is because if we show inequality (\ref{ephi1}) for the case where $\phi=0$, then by making a variable change from $H_1,\dots, H_{r-1}$ to $H_j^\prime\coloneqq e^{\phi} H_j\ (j=1,\dots, r-1)$, we have inequality (\ref{ephi1}) for a general $\phi$ and for the same $C$ as in the case where $\phi=0$. We see that the above variable transformation is consistent since we can check the following:
\begin{align*}
H_r^\prime=e^\phi H_r,
\end{align*} 
where $H_r^\prime$ is defined to be $H_r^\prime\coloneqq H_1^{\prime-1}\cdot H_2^{\prime-1}\cdots H_{r-1}^{\prime-1}\cdot e^{r\phi}$. We show inequality (\ref{ephi1}) for the case where $\phi=0$. We arrange the method given in \cite[Proof of Proposition 3.8]{LM1} for our proof. We set $c_\max\coloneqq \max\{H_j\mid 1\leq j\leq r\}$ and $c_\min\coloneqq \min \{H_j\mid 1\leq j\leq r\}$. We take $j_\max$ and $j_\min$ so that $c_\max=H_{j_\max}$ and $c_\min=H_{j_\min}$. 
In what follows, for simplicity we will only consider the case where they are 1 and $r$, respectively. Even if this were not the case, a simple consideration shows that the following estimates would hold just as well, with a slight modification of the argument. The following holds:
\begin{align*}
c_\max-c_\min&=\sum_{j=2}^r(H_{j-1}-H_j)\\
&\leq (\sum_{j=2}^r(H_{j-1}-H_j)^2)^{1/2}(\sum_{j=2}^r1^2)^{1/2}\\
&\leq \sqrt{r-1}(\sum_{j=1}^r(H_{j-1}-H_j)^2)^{1/2},
\end{align*}
where we have used the Cauchy-Schwartz inequality. Therefore we have 
\begin{align}
\sum_{j=1}^r(H_{j-1}-H_j)^2&\geq \frac{1}{r-1}(c_\max-c_\min)^2 \notag \\
&=\frac{1}{r-1}c_\max^2(1-c_\min/c_\max)^2. \label{cmax}
\end{align}
On the other hand, we have
\begin{align*}
\sum_{j=1}^rH_j-r\leq rc_\max-r.
\end{align*}
Since $\Pi_{i=1}^rH_i=H_1\cdot H_2\cdots H_r=1$, it holds that $\sum_{j=1}^rH_j-r\geq 0$. Hence we have
\begin{align}
(\sum_{j=1}^rH_j-r)^2&\leq r^2(c_\max-1)^2 \notag \\
&=r^2c_\max^2(1-1/c_\max)^2 \label{cmin}.
\end{align}
We compare (\ref{cmax}) and (\ref{cmin}). Since $c_\min\leq 1$, we have
\begin{align*}
1-c_\min/c_\max\geq 1-1/c_\max.
\end{align*}
Furthermore, since $1-1/c_\max\geq 0$, we have
\begin{align}
(1-c_\min/c_\max)^2\geq (1-1/c_\max)^2. \label{cmaxmin}
\end{align}
Then from (\ref{cmax}), (\ref{cmin}), and (\ref{cmaxmin}), for the case where $\phi\neq -\infty$, we have
\begin{align}
\sum_{j=1}^r(H_{j-1}-H_j)^2\geq \frac{1}{r^2(r-1)}(\sum_{j=1}^rH_j-re^\phi)^2. \label{rsquare}
\end{align}
We next consider the case where $\phi=-\infty$, i.e., $e^\phi=0$. Let $u_1,\dots, u_r$ be the canonical basis of $\R^{r-1}$. We define invertible $(r-1)\times (r-1)$-matrices $A$ and $B$ as follows:
\begin{align*}
A u_1&\coloneqq u_1, \\
A u_j&\coloneqq u_j-u_{j-1} \ (j=2,\dots, r-1), 
\end{align*}
where $u_1,\dots, u_{r-1}$ are considered to be column vectors. We set $\hat{H}\coloneqq (H_1,\dots, H_{r-1})\in\R^{r}$. Then we have
\begin{align*}
\sum_{j=1}^r(H_{j-1}-H_j)^2&=H_1^2+\sum_{j=2}^{r-1}(H_{j-1}-H_j)^2+H_{r-1}^2 \\
&\geq H_1^2+\sum_{j=2}^{r-1}(H_{j-1}-H_j)^2 \\
&=|A\hat{H}|^2,
\end{align*}
where the norm $|\cdot|$ denotes the Euclidean norm. On the other hand, by the Cauchy-Schwartz inequality, we have
\begin{align*}
(\sum_{j=1}^{r-1}H_j)^2&\leq (r-1)|\hat{H}|^2 \\
&\leq (r-1)|A^{-1}\cdot A\hat{H}|^2 \\
&\leq (r-1)|A^{-1}|^2|A\hat{H}|^2 \\
&=(r-1)\cdot \frac{1}{2}r(r-1) |A\hat{H}|^2\\
&=\frac{1}{2}r(r-1)^2|A\hat{H}|^2.
\end{align*}
Therefore we have
\begin{align}
\sum_{j=1}^r(H_{j-1}-H_j)^2\geq \frac{2}{r^2(r-1)}(\sum_{j=1}^{r-1}H_j^2). \label{rsquare2}
\end{align}
for the case where $\phi=-\infty$. From inequalities (\ref{rsquare}) and (\ref{rsquare2}), we have inequality (\ref{ephi1}) with $C=\frac{1}{r^2(r-1)}$.
\end{proof}
\begin{rem} 
This is a digression and not directly related to the main topic, but when $r=2$, the constant $C$ in inequality (\ref{ephi1}) can be taken to be $2$. The author does not know the optimal estimate of inequality (\ref{ephi1}).
\end{rem}
\begin{lemm}[cf. \cite{LM1, Sim2}]\label{lem2} {\it Let $C$ be a positive constant as in Lemma \ref{lem1}. Suppose that there exists a positive real number $H>0$ such that $e^\phi\leq H$ and that $\sum_{j=1}^rH_j\geq 2r H$. Then the following holds:
\begin{align*}
\sum_{j=1}^r(H_{j-1}-H_j)^2\geq \frac{C}{4}(\sum_{j=1}^rH_j)^2.
\end{align*}
}
\end{lemm}
\begin{proof} The proof is the same as in the proof of \cite[Lemma 3.6]{LM1}. From the assumption, we have
\begin{align*}
\sum_{j=1}^rH_j-re^\phi\geq rH \geq re^\phi.
\end{align*}
This implies
\begin{align*}
\frac{1}{2}\sum_{j=1}^r H_j \geq re^\phi,
\end{align*}
and thus 
\begin{align*}
\sum_{j=1}^rH_j-re^\phi\geq \frac{1}{2}\sum_{j=1}^r H_j.
\end{align*}
From this we have the desired inequality.
\end{proof}

For the proof of the following Lemma \ref{lem3} and Lemma \ref{lem4}, we refer the reader to \cite[Proposition 3.8]{LM1} and \cite[Lemma 3.10]{LM1}, respectively.
\begin{lemm}[\cite{LM1}]\label{lem3}{\it For any $0<\epsilon<1$, there exists $\delta>0$ that depends only on $\epsilon$ and $r$, such that the following holds:
\begin{itemize}
\item If there exists $1\leq j\leq r-1$ such that $H_j\leq \epsilon e^\phi$, then we obtain $\sum_{j=1}^r(H_{j-1}-H_j)^2\geq \delta (\sum_{j=1}^r H_j)^2$.
\end{itemize}
}
\end{lemm}
Let $h_1^\prime,\dots, h_r^\prime$ be another positive real numbers satisfying $\Pi_{i=1}^rh_i^\prime=h_1^\prime\cdot h_2^\prime\cdots h_{r-1}^\prime\cdot h_r^\prime=1$. We set $s_j\coloneqq h_j^{-1}\cdot h_j^\prime$ for each $j=1,\dots, r$. We also set $s\coloneqq (s_1,\dots, s_r)\in (\R_{>0})^r$. We assume that there exist positive constants $B$ and $C$ such that the following conditions hold:
\begin{itemize}
\item $H_j\leq B$ for each $j=1,\dots, r-1$.
\item $\sum_{j=1}^rs_j\leq C$.
\end{itemize}
Then the following lemma is quoted from \cite[Lemma 3.10]{LM1}:
\begin{lemm}[\cite{LM1}]\label{lem4}{\it There exists $C_1>0$ and $\epsilon_1>0$, depending only on $B$ and $C$, such that the following holds:
\begin{itemize}
\item If there exists $0<\epsilon<\epsilon_1$ such that $\sum_{j=1}^rs_j^{-1}(s_{j+1}-s_j)^2H_j \leq \epsilon^2$, then it holds that $\sum_{j=1}^r|s_j-1|\leq C_1\cdot \epsilon$.
\end{itemize}
}
\end{lemm}
\begin{rem} The consistency between the notation $|[f,s]s^{-1/2}|_h$ used in \cite[Lemma 3.10]{LM1} and the one we used in the above is as follows (see also \cite[Lemma 3.1]{Sim1} and \cite[Proof of Lemma 3.10]{LM1}): Let $h=(h_1,\dots, h_r)$ and $h^\prime=(h_1^\prime,\dots, h_r^{\prime})$ be diagonal Hermitian metrics on $(\K_r,\Phi(q))\rightarrow X$. We set $s_j\coloneqq h_j^{-1}\otimes h_j^\prime$ for each $j=1,\dots, r$, $s_0=s_r$, and $s\coloneqq (s_1,\dots, s_r)$. We also define Hermitian metrics $H_1,\dots, H_{r-1}, H_r=H_0$ on $K_X^{-1}\rightarrow X$ in the same way as in Section \ref{2.1}. Let $u_1$ be a local frame around $x\in X$ of $K_X^{(r-1)/2}\rightarrow X$. Let $z$ be a local holomorphic chart around $x$ and we define a local frame $u_j$ of $K_X^{\frac{r-(2j-1)}{2}}\rightarrow X$ as $u_j\coloneqq u_1\otimes (\frac{\partial}{\partial z})^{(j-1)}$ for each $j=2,\dots, r-1$. Then we can check that $\Phi(q)^{\ast h}u_{j+1}=H_j(\frac{\partial}{\partial z}, \frac{\partial}{\partial z})u_j\otimes d\bar{z}$ and the following:
\begin{align*}
&\inum[\Phi(q),s]s^{-1/2}\wedge ([\Phi(q),s]s^{-1/2})^{\ast h} u_j \\
=&\inum[\Phi(q),s]s^{-1}[s,\Phi(q)^{\ast h}]u_j \\
=&s_{j-1}^{-1}(s_j-s_{j-1})^2H_{j-1}\left(\frac{\partial}{\partial z},\frac{\partial}{\partial z}\right)u_j\inum dz\wedge d\bar{z}
\end{align*}
for each $j=1,\dots,r$. From this, we can verify the consistency between the notation used in \cite{LM1} and our notation.
\end{rem}
Following \cite[Section 3]{LM1}, we apply the above lemmas on linear algebras to a priori estimates of solution to equation (\ref{phi}). Let $h=(h_1,\dots, h_r)$ be a weak solution to equation (\ref{phi}) satisfying $\det(h)=1$. As stated in \cite{Miy3}, it follows from the general regularity theory of the Poisson equation that $h$ is at least of class $C^{1,\alpha}$ for any $\alpha\in (0,1)$. We denote by $H_1,\dots, H_r$ the Hermitian metrics on $K_X^{-1}\rightarrow X$ induced by $h$ as in Section \ref{2.1}.
The following is an extension of the well-known inequality for Higgs bundles established by Simpson \cite[Lemma 10.1]{Sim1} (for the proof, see \cite[Proposition 11]{Miy3} and \cite[Proposition 23]{Miy5}):
\begin{lemm}\label{free energy}{\it The following inequality holds in the sense of the distribution: 
\begin{align}
\inum\partial\bar{\partial} \log(\sum_{j=1}^r\vol(H_j)/\vol(H_\refe)) \geq \frac{\sum_{j=1}^r(\vol(H_{j-1})-\vol(H_j))^2}{\sum_{j=1}^r\vol(H_j)}+\inum F_{H_\refe}. \label{well-known}
\end{align}
}
\end{lemm}
By using the above Lemma \ref{free energy}, we can show the following as in \cite[Lemma 3.13]{LM1}:
\begin{lemm}\label{ephi}{\it Suppose that there exists a Hermitian metric $H$ on $K_X^{-1}\rightarrow X$ such that
\begin{align*}
e^{\varphi}h_\refe^{-1}\leq H.
\end{align*}
Then there exist positive constants $C_1$ and $C_2$ independent of $H$ such that
\begin{align}
\inum\Lambda_H\partial\bar{\partial}\log(\sum_{j=1}^r\vol(H_j)/\vol(H))\geq C_1(\sum_{j=1}^r\vol(H_j)/\vol(H))-C_2+\inum \Lambda_HF_H. \label{C1C2}
\end{align}
where we denote by $\Lambda_H=(\vol(H)\wedge)^{\ast}$ the contraction operator induced by $H$.
}
\end{lemm}
\begin{proof} In what follows, we will identify the Hermitian metrics on $K_X^{-1}\rightarrow X$ with positive real numbers at each point by fixing one trivialization for each point. Note, however, that all the following discussions are independent of the choice of trivialization. From Lemma \ref{lem1}, there exists a positive constant $C$ depending only on $r$ such that
\begin{align*}
\sqrt{\sum_{j=1}^r(H_{j-1}-H_j)^2}\geq \sqrt{C}(\sum_{j=1}^rH_j-re^{\varphi}h_\refe^{-1}).
\end{align*}
From the assumption $e^\varphi h_\refe^{-1}\leq H$, we have
\begin{align}
\sqrt{\sum_{j=1}^r(H_{j-1}-H_j)^2}\geq \sqrt{C}(\sum_{j=1}^rH_j-re^{\varphi}h_\refe^{-1})\geq \sqrt{C}(\sum_{j=1}^rH_j-2r H). \label{sqrtc1}
\end{align}
On the other hand, from Lemma \ref{lem2}, under the assumption of $\sum_{j=1}^rH_j\geq 2rH$, we have
\begin{align}
\sqrt{\sum_{j=1}^r(H_{j-1}-H_j)^2}\geq \frac{\sqrt{C}}{2}\sum_{j=1}^rH_j. \label{sqrtc2}
\end{align}
From inequalities (\ref{sqrtc1}) and (\ref{sqrtc2}), we have
\begin{align}
\sum_{j=1}^r(H_{j-1}-H_j)^2\geq \frac{C}{2}(\sum_{j=1}^rH_j)(\sum_{j=1}^rH_j-2rH).
\end{align}
Therefore, the following holds globally over $X$:
\begin{align}
\frac{\sum_{j=1}^r(\vol(H_{j-1})-\vol(H_j))^2}{\sum_{j=1}^r\vol(H_j)}\geq \frac{C}{2}(\sum_{j=1}^r\vol(H_j)-2r\vol(H)). \label{-2rH}
\end{align}
From the above inequality (\ref{-2rH}) and inequality (\ref{well-known}), we have the desired conclusion.
\end{proof}
The following Corollary \ref{CYM} is an extension of the crucial insight discovered by Li-Mochizuki \cite[Section 3.3]{LM1} to more general subharmonic weight functions:
\begin{cor}\label{CYM}{\it Suppose that Condition \ref{condition} is satisfied. Let $g$ be a complete K\"ahler metric on $X$ whose Ricci curvature is bounded below and $H$ the Hermitian metric on $K_X^{-1}\rightarrow X$ induced by $g$. The for the K\"ahler metric $g$, the following (i) and (ii) are equivalent:
\begin{enumerate}[(i)]
\item The function $e^{\varphi}h_\refe^{-1}\otimes H^{-1}$ is a bounded function.
\item For each $j=1,\dots, r$, $H_j\otimes H^{-1}$ is a bounded function. 
\end{enumerate}
}
\end{cor}
\begin{proof} It is clear that (ii) implies (i). We prove the converse. We suppose that (i) holds. For simplicity, we assume that we can take 1 as an upper bound on the function $e^\varphi h_\refe^{-1}\otimes H^{-1}$. Let $Y\subseteq X$ be an open subset with smooth boundary such that $K\subseteq Y$ and that the closure $\overline{Y}$ is compact. 
Note that since Condition \ref{condition} is satisfied, $\log(\sum_{j=1}^r\vol(H_j)/\vol(H))$ is a $C^2$ function on $X\backslash \overline{Y}$. By applying the Cheng-Yau maximum principle for functions on manifolds with boundary to inequality (\ref{C1C2}) on $X\backslash Y$, it holds that $\log(\sum_{j=1}^r\vol(H_j)/\vol(H))$ is a bounded function on $X\backslash Y$ and therefore on $X$. This implies the assertion of (ii).
\end{proof}
Throughout this section, we adapt the notation denoting by $\Lambda_H$ the contraction operator induced by a Hermitian metric $H$ on $K_X^{-1}\rightarrow X$. Following \cite[Lemma 3.22]{LM1}, we can show the following:
\begin{lemm}\label{log}{\it There exist positive constants $C_1$ and $C_2$ such that the following holds for each $i=1,\dots, r-1$:
\begin{align}
\inum\Lambda_{H_i}\partial\bar{\partial} \log(\sum_{j=1}^r(\vol(H_j)/\vol(H_i))) \geq C_1(\sum_{j=1}^r\vol(H_j)/\vol(H_i))-C_2+\inum \Lambda_{H_i}F_{H_i}. \label{Hi}
\end{align}
}
\end{lemm}
\begin{proof} As in the proof of Lemma \ref{ephi}, we identify the Hermitian metrics on $K_X^{-1}\rightarrow X$ with positive real numbers at each point by fixing one trivialization for each point. For each $i=1,\dots, r-1$, if $2H_i\geq e^{\varphi}h_\refe^{-1}$, then by the same argument as in the proof of Lemma \ref{ephi}, we have
\begin{align*}
\frac{\sum_{j=1}^r(H_{j-1}-H_j)^2}{\sum_{j=1}^rH_j}\geq \frac{C}{2}(\sum_{j=1}^rH_j-4rH_i),
\end{align*}
where $C$ is a positive constant as in Lemma \ref{lem1} and Lemma \ref{lem2}. On the other hand, if $2H_i<e^{\varphi}h_\refe^{-1}$, then from Lemma \ref{lem3}, there exists a positive constant $\delta$ depending only on $r$ (and the choice of $\epsilon=\frac{1}{2}$) such that
\begin{align*}
\frac{\sum_{j=1}^r(H_{j-1}-H_j)^2}{\sum_{j=1}^rH_j}\geq \delta (\sum_{j=1}^rH_j).
\end{align*}
From the above argument, the following holds globally on $X$:
\begin{align}
\frac{\sum_{j=1}^r(\vol(H_{j-1})-\vol(H_j))^2}{\sum_{j=1}^r\vol(H_j)}\geq \min\{C/2,\delta\}(\sum_{j=1}^r\vol(H_j))-2rC\vol(H_i). \label{c1}
\end{align}
From inequality (\ref{c1}) and inequality (\ref{well-known}), we have the desired conclusion.
\end{proof}
\begin{cor}\label{ji}{\it Suppose that Condition \ref{condition} is satisfied and that an $H_i$ among $H_1, \dots, H_{r-1}$ is complete. Then $H_j\otimes H_i^{-1}$ is a bounded function for all $j=1,\dots, r$. 
}
\end{cor}
\begin{proof} Let $Y\subseteq X$ be an open subset with smooth boundary such that $K\subseteq Y$ and that the closure $\overline{Y}$ is compact. By the same reason as in \cite[Lemma 2.4]{LM1}, the Gaussian curvature of $H_i$ is bounded below. By applying the Cheng-Yau maximum principle for functions on mainifolds with boundary to inequality (\ref{Hi}) on $X\backslash Y$, $\log(\sum_{j=1}^r\vol(H_j)/\vol(H_i))$ is a bounded function on $X\backslash Y$ and therefore on $X$. This implies the claim.
\end{proof}
Let $h^\prime=(h_1^\prime,\dots, h_r^\prime)$ be another solution to equation (\ref{phi}) satisfying $\det(h^\prime)=1$. We denote by $s=(s_1,\dots, s_r)$ the endomorphism determined by $h(s\cdot,\cdot)=h^\prime(\cdot,\cdot)$. We set $s_0\coloneqq s_r$.
\begin{lemm}\label{ddbars}{\it The endomorphism $s$ satisfies the following inequality in the sense of the distribution:
\begin{align}
\inum\partial\bar{\partial} \sum_{j=1}^rs_j\geq \sum_{j=1}^rs_{j-1}^{-1}(s_{j-1}-s_j)^2\vol(H_{j-1}). \label{s}
\end{align}
}
\end{lemm}
\begin{proof} By a direct computation, we have
\begin{align*}
\inum\partial \bar{\partial} \sum_{j=1}^rs_j&= \inum\partial \bar{\partial}\sum_{j=1}^re^{\log s_j} \\
&= \inum \partial \sum_{j=1}^r(\bar{\partial}\log s_j) e^{\log s_j} \\
&= \inum \sum_{j=1}^r\partial (\log s_j)\wedge\bar{\partial}(\log s_j) e^{\log s_j}+\inum\sum_{j=1}^r(\partial \bar{\partial} \log s_j )e^{\log s_j} \\
&\geq \inum \sum_{j=1}^rs_j\partial \bar{\partial}(\log s_j) \\
&=\sum_{j=1}^rs_j(\inum F_{h_j}-\inum F_{h_j^\prime}) \\
&=\sum_{j=1}^rs_j(-\vol(H_{j-1})+\vol(H_j)+\vol(H_{j-1}^\prime)-\vol(H_j^\prime)). 
\end{align*}
On the other hand, since we have $s_j\cdot s_{j-1}^{-1}=H_{j-1}^{-1}\otimes H_{j-1}^\prime$ for each $j=1,\dots, r$, it holds that
\begin{align*}
s_{j-1}^{-1}(s_j-s_{j-1})^2\vol(H_{j-1})&=(s_j-s_{j-1})(s_js_{j-1}^{-1}-1)\vol(H_{j-1}) \\
&=(s_j-s_{j-1})(\vol(H_{j-1}^\prime)-\vol(H_{j-1})).
\end{align*}
Therefore, we have
\begin{align*}
\sum_{j=1}^rs_{j-1}^{-1}(s_{j-1}-s_j)^2\vol(H_{j-1}) 
&=\sum_{j=1}^r(s_j-s_{j-1})(\vol(H_{j-1}^\prime)-\vol(H_{j-1})) \\
&=\sum_{j=1}^r(s_j\vol(H_{j-1}^\prime)+s_{j-1}\vol(H_{j-1})-s_{j-1}\vol(H_{j-1}^\prime)-s_j\vol (H_{j-1})) \\
&=\sum_{j=1}^rs_j(\vol(H_{j-1}^\prime)+\vol(H_j)-\vol(H_j^\prime)-\vol(H_{j-1})). \\
&=\sum_{j=1}^rs_j(-\vol(H_{j-1})+\vol(H_j)+\vol(H_{j-1}^\prime)-\vol(H_j^\prime)). 
\end{align*}
This implies the claim.
\end{proof}
Then we prove Theorem \ref{main theorem 1}.
\begin{proof}[Proof of Theorem \ref{main theorem 1}] By Corollary \ref{CYM} and Corollary \ref{ji}, under the assumption of Condition \ref{condition}, $H_1,\dots, H_{r-1}$ and $H_1^\prime,\dots, H_{r-1}^\prime$ are all mutually bounded. Therefore $s$ is a bounded function. We choose any one of $H_1,\dots, H_{r-1}, H_1^\prime, \dots, H_{r-1}^\prime$ and write it as $H$. By the same reason as \cite[Lemma 2.4]{LM1}, the Gaussian curvature of $H$ is bounded below. From inequality (\ref{s}), we have the following inequality 
\begin{align}
\inum\Lambda_H\partial\bar{\partial} \sum_{j=1}^rs_j\geq \sum_{j=1}^rs_{j-1}^{-1}(s_{j-1}-s_j)^2H_{j-1}\otimes H^{-1}. \label{sH}
\end{align}
Let $Y\subseteq X$ be an open subset with smooth boundary such that $K\subseteq Y$ and that the closure $\overline{Y}$ is compact. By applying the Omori-Yau maximum principle for functions on manifolds with boundary to inequality (\ref{sH}), one of the following holds:
\begin{enumerate}
\item It holds that $\sup_{X\setminus Y}s=\max_{\partial Y}s$, where $\partial Y$ is the boundary $\overline{Y}\setminus Y$of $Y$.
\item There exists $m_0\in\Z_{\geq 1}$ and a sequence of points $(x_m)_{m\geq m_0}$ on $X\setminus Y$ such that
\begin{align*}
&(\sum_{j=1}^rs_j)(x_m)\geq \sup_{X\setminus Y}(\sum_{j=1}^rs_j)-1/m, \\
&(\inum\Lambda_H\partial \bar{\partial} \sum_{j=1}^rs_j)(x_m)\leq 1/m.
\end{align*}
\end{enumerate}
We first consider the case 2. From inequality (\ref{sH}), we see that $s=(s_1,\dots, s_r)$ satisfies the following inequality at each $x_m$:
\begin{align*}
\sum_{j=1}^r(s_{j-1}^{-1}(s_{j-1}-s_j)^2)(x_m)(H_{j-1}\otimes H^{-1})(x_m)\leq 1/m.
\end{align*}
Around each $x_m\ (m\geq m_0)$, we fix a coordinate such that $H(\frac{\partial}{\partial z}, \frac{\partial }{\partial z})(x_m)=1$. 
Since $H_j\otimes H^{-1}$ is a bounded function for all $j=1,\dots, r$, there exists $B>0$ such that $H_j(\frac{\partial}{\partial z}, \frac{\partial}{\partial z})(x_m)\leq B$ for all $j=1,\dots, r$ and for all $m\geq m_0$. Then from Lemma \ref{lem4}, there exists $m_1\geq m_0$ and $C_1>0$ such that 
\begin{align*}
\sum_{j=1}^r|s_j(x_m)-1|\leq C_1/m^{1/2}
\end{align*}
for all $m\geq m_1$. This implies that $\sup_{X\setminus Y}\sum_{j=1}^rs_j\leq r$. Since we have $s_1\cdots s_r=1$, it holds that $\sum_{j=1}^rs_j\geq r$. As a conclusion, we have $\sum_{j=1}^r s_j = r$ on $X \setminus Y$.
Therefore, in both Case 1 and Case 2, we see that $\sum_{j=1}^r s_j$ attains its maximum value somewhere on $\overline{Y}$, viewed as a function on $X$. From inequality (\ref{s}), we see that $\sum_{j=1}^r s_j$ is a subharmonic function. Therefore, by the maximum principle for subharmonic functions, we find that $\sum_{j=1}^r s_j$ is constant.
By repeating the argument of Case 1 for $s$ in Case 2, we find that $\sum_{j=1}^r s_j = r$ holds in both cases. Therefore, we have $h = h'$.
\end{proof}

\subsection{Supersolution and subsolution to equation (\ref{phi})}
Before proving Theorems \ref{main theorem 2} and \ref{main theorem 3}, this subsection prepares some lemmas on supersolutions and subsolutions to equation (\ref{phi}). Let $e^{-\varphi}h_\refe$ a semipositive singular Hermitian metric on $K_X\rightarrow X$ with a smooth reference metric $h_\refe$. We denote by $(h_{\refe,1},\dots, h_{\refe,r})$ the real diagonal Hermitian metric on $\K_r\rightarrow X$ induced by $h_\refe$. Let $h = (h_1, \dots, h_r)=(e^{-\phi_1}h_{\refe,1}, e^{-\phi_2}h_{\refe,2}, \dots, e^{-\phi_r}h_{\refe,r})$ be a real diagonal Hermitian metric on $\K_r \rightarrow X$. We set $n\coloneqq [r/2]$. Suppose that for each $j=1,\dots,n$, the weight function $\phi_j$ is $L^1_{loc}$ and for each $j=0,\dots, n$, $\phi_j-\phi_{j+1}$ is $L^\infty_{loc}$, where $\phi_0$ is $\phi_r$. We set $H_j\coloneqq h_j^{-1}\otimes h_{j+1}$ for each $j=1,\dots, r-1$ and $H_0=H_r\coloneqq e^{r\varphi}h_\refe^{-1}\otimes h_r^{-1}\otimes h_1$. The notion of supersolution and subsolution to equation (\ref{phi}) is then defined as follows (cf. \cite[Definition 5.1]{LM1}):
\begin{defi} We say that $h$ is a {\it supersolution} (resp. {\it subsolution}) to equation \eqref{phi} if it satisfies the following differential inequality in the weak sense:
\begin{align*}
\inum F_{h_j} + \vol(H_{j-1}) - \vol(H_j) \geq 0 \quad (\text{resp. } \leq 0) \quad \text{for all } j = 1, \dots, n.
\end{align*}
\end{defi}
\begin{ex} Let $h=(h_1,\dots, h_r)$ be the diagonal Hermitian metric on $\K_r\rightarrow X$ induced by $e^{-\varphi}h_\refe$. Then we have $H_1=\cdots= H_r=e^{\varphi}h_\refe^{-1}$. Therefore $h$ is a supersolution to equation (\ref{phi}).
\end{ex}
\begin{ex} Let $h_+=(h_{+,1},\dots, h_{+,r})$ and $h_+^\prime=(h_{+,1}^\prime, \dots, h_{+,r}^\prime)$ be two continuous supersolutions to equation (\ref{phi}), and let $h_-=(h_{-,1},\dots, h_{-,r})$ and $h_-^\prime=(h_{-,1}^\prime, \dots, h_{-,r}^\prime)$ be two continuous subsolutions to equation (\ref{phi}). For each $j=1,\dots, n$, we define $\widetilde{h}_{+,j}$ and $\widetilde{h}_{-,j}$ as follows:
\begin{align*} 
\widetilde{h}_{+,j}&\coloneqq \min\{h_{+,j}, h_{+,j}^\prime\}, \\
\widetilde{h}_{-,j}&\coloneqq \max\{h_{-,j}, h_{-,j}^\prime\}. 
\end{align*}
For the remaining indices $j=n+1,\dots, r$, we define $\widetilde{h}_{+,j}$ and $\widetilde{h}_{-,j}$ so that $\widetilde{h}_{+}\coloneqq (\widetilde{h}_{+,1},\dots, \widetilde{h}_{+,r})$ and $\widetilde{h}_{-}\coloneqq (\widetilde{h}_{-,1},\dots, \widetilde{h}_{-,r})$ are real. Then the Hermitian metric $\widetilde{h}_+$ (resp. $\widetilde{h}_-$) is a supersolution (resp. subsolution) to equation (\ref{phi}) (cf. \cite[Lemma 5.4]{LM1}).
\end{ex}
\begin{ex}\label{subandsuper} Let $\varphi$ and $\varphi^\prime$ be subharmonic weight functions. Let $h_+$ and $h_-$ be a real diagonal supersolution and subsolution to equation (\ref{phi}) associated with $\varphi$, respectively. Suppose that $\varphi \leq \varphi^\prime$. Then $h_+$ is also a supersolution to equation (\ref{phi}) associated with $\varphi^\prime$. Similarly, if $\varphi \geq \varphi^\prime$, then $h_-$ is also a subsolution to equation (\ref{phi}) associated with $\varphi^\prime$.
\end{ex}
Following \cite[Proposition 4.6]{LM1}, we prepare the following lemma:
\begin{lemm}\label{cands}{\it Suppose that $e^\varphi$ is smooth. Let $h=(h_1,\dots, h_r)$ be a smooth complete subsolution to equation (\ref{phi}) and $h^\prime=(h_1^\prime,\dots, h_r^\prime)$ a smooth supersolution to equation (\ref{phi}). Suppose that $h_j \otimes h_j^{\prime-1}$ is a bounded function for each $j=1,\dots, n$. Then one of the following holds: 
\begin{enumerate}[(i)]
\item $h_j\otimes h_j^{\prime-1}<1\ \text{for all $j=1,\dots, n$}$; 
\item $h_j=h_j^\prime$ for all $j=1,\dots, n$.
\end{enumerate}
}
\end{lemm}
\begin{proof}
For each $j=1,\dots, n$, we define $\xi_j\coloneqq \log(h_j\otimes h_j^{\prime-1})$ and $M_j\coloneqq \sup_Xe^{\xi_j}$. We set $H_j^\prime\coloneqq h_j^{\prime-1}\otimes h_{j+1}^\prime$ for each $j=1,\dots, r-1$ and $H_0^\prime=H_r^\prime\coloneqq h_r^{\prime-1}\otimes h_1^\prime\otimes (e^{-\varphi}h_\refe)^{-r}$. Suppose that (ii) does not hold, i.e., there exists at least one $j = 1, \dots, n$ such that $\xi_j$ is not constant. Following \cite[Proof of Proposition 4.6]{LM1}, we show that $e^{\xi_j}<1$ for each $j=1,\dots, n$. By direct calculation, we have the following:
\begin{align}
\inum\partial \bar{\partial} \xi_j\geq \vol(H_{j-1})-\vol(H_j)-\vol(H_{j-1}^\prime)+\vol(H_j^\prime) \ \text{for $j=1,\dots, n$}. \label{xij}
\end{align}
Suppose that $n=1$. Then from (\ref{xij}), we have
\begin{align}
\inum \Lambda_{H_1}\partial \bar{\partial} \xi_1&\geq (1-e^{-2\xi_1})H_0\otimes H_1^{-1}-(1-e^{(4-r)\xi_1}). \label{xi01}
\end{align}
Note that by the same reason as in \cite[Lemma 2.4]{LM1}, the Gaussian curvature of $H_1$ is bounded below and that $\xi_1$ is bounded above from the assumption. Therefore, we can apply the Omori-Yau maximum principle to inequality (\ref{xi01}). By the Omori-Yau maximum principle, there exist $m_0\in\Z_{\geq 1}$ and a sequence of points $(x_m)_{m\geq m_0}$ on $X$ such that
\begin{align}
\xi_1(x_m)&\geq M_1-1/m, \label{mmm}\\
1/m&\geq (1-e^{-2\xi_1(x_m)})(H_0\otimes H_1^{-1})(x_m)-(1-e^{(4-r)\xi_1(x_m)}). \label{1mxm}
\end{align}
From (\ref{mmm}) and (\ref{1mxm}), we have $1\geq M_1$. Then since $\xi_1$ is not constant, inequality (\ref{xi01}) and the strong maximum principle immediately yields $M_1<1$. We next consider the case where $n\geq 2$. From (\ref{xij}), we have
\begin{align}
\inum\Lambda_{H_1}\partial \bar{\partial} \xi_1&\geq (1-e^{-2\xi_1})H_0\otimes H_1^{-1}-(1-e^{\xi_1-\xi_2}), \label{xi12}\\
\inum\Lambda_{H_j} \partial \bar{\partial} \xi_j &\geq (1-e^{\xi_{j-1}-\xi_j})H_{j-1}\otimes H_j^{-1}-(1-e^{\xi_j-\xi_{j+1}}) \ \text{for $j=2,\dots, n-1$}, \label{xi34}\\
\inum \Lambda_{H_n}\partial \bar{\partial} \xi_n&\geq (1-e^{\xi_{n-1}-\xi_n})H_{n-1}\otimes H_n^{-1}-(1-e^{(2n+2-r)\xi_n}). \label{xin}
\end{align}
By the same reason as in \cite[Lemma 2.4]{LM1}, the Gaussian curvature of $H_i$ is bounded below for each $i=1,\dots, n$. Also, from the assumption, $\xi_1,\dots, \xi_n$ are bounded above. Therefore, we can apply the Omori-Yau maximum principle to inequalities (\ref{xi12}), (\ref{xi34}), and (\ref{xin}). 
From (\ref{xi12}) and the Omori-Yau maximum principle, there exist $m_0\in\Z_{\geq 1}$ and a sequence of points $(x_m)_{m\geq m_0}$ on $X$ such that
\begin{align}
\xi_1(x_m)&\geq M_1-1/m, \label{mmm2}\\
\frac{1}{m}&\geq (1-e^{-2\xi_1(x_m)})(H_0\otimes H_1^{-1})(x_m)-(1-e^{\xi_1(x_m)-\xi_2(x_m)}) \notag \\
&\geq (1-e^{-2\xi_1(x_m)})(H_0\otimes H_1^{-1})(x_m)-(1-M_2^{-1}e^{\xi_1(x_m)}). \label{1mxm2}
\end{align}
From (\ref{mmm2}) and (\ref{1mxm2}), if $M_1\geq 1$, then we have $M_1\leq M_2$. Therefore, we obtain either $M_1<1$ or $M_1\leq M_2$. By the same argument using the Omori-Yau maximum principle as above, from (\ref{xi34}), we obtain either $M_j< M_{j-1}$ or $M_j\leq M_{j+1}$ for each $j=2,\dots, n-1$ and from (\ref{xin}) we obtain either $M_n<M_{n-1}$ or $M_n\leq 1$. We define a subset $A$ of $\{1,\dots, n\}$ as $A\coloneqq \{j\mid M_j=\max_{1\leq k\leq n}M_k\}$. If $j\in A$ for some $2\leq j\leq n-1$, then we have $M_j \leq M_{j+1}$, and hence $j+1\in A$, since $M_j < M_{j-1}$ does not hold. By repeating this argument, we obtain $k\in A$ for all $j\leq k\leq n$ and $M_n\leq 1$. If $1\in A$, then we obtain either $M_1<1$ or $M_1=M_2$. In the latter case, from the discussion already given, we see that $j\in A$ for all $1\leq j\leq n$ and $M_n\leq 1$. If $n\in A$, we have $M_n\leq 1$. As a conclusion, in either case, we have $M_j\leq 1$ for all $1\leq j\leq n$. From (\ref{xi12}), (\ref{xi34}), and (\ref{xin}), we have
\begin{align}
\inum\Lambda_{H_1}\partial \bar{\partial} \xi_1&\geq (1-e^{-2\xi_1})H_0\otimes H_1^{-1}-(1-e^{\xi_1}), \label{xi12+}\\
\inum\Lambda_{H_j} \partial \bar{\partial} \xi_j &\geq (1-e^{-\xi_j})H_{j-1}\otimes H_j^{-1}-(1-e^{\xi_j}) \ \text{for $j=2,\dots, n-1$}, \label{xi34+}\\
\inum \Lambda_{H_n}\partial \bar{\partial} \xi_n&\geq (1-e^{-\xi_n})H_{n-1}\otimes H_n^{-1}-(1-e^{(2n+2-r)\xi_n}). \label{xin+}
\end{align}
By the strong maximum principle, we obtain either $\xi_j<0$ or $\xi_j=0$ for each $j=1,\dots, n$. If $\xi_1=0$, then from (\ref{xi12}) and $M_2\leq 1$, we have $\xi_2=0$. If $\xi_j=0$ for some $j=2,\dots, n-1$, then from (\ref{xi34}), $M_{j-1}\leq 1$, and $M_{j+1}\leq 1$, we have $\xi_{j-1}=\xi_{j+1}=0$. If $\xi_n=0$, then from (\ref{xin}) and $M_{n-1}\leq 1$, we have $\xi_{n-1}$. Since we assume that there exists a $j=1,\dots, n$ such that $\xi_j$ is not constant, we conclude that $e^{\xi_j}<1$ for all $j=1,\dots, n$.
\end{proof}



\subsection{Proof of Theorems \ref{main theorem 2} and \ref{main theorem 3}}
This subsection proves Theorems \ref{main theorem 2} and \ref{main theorem 3}. Suppose that $X$ is the unit disc $\D\coloneqq \{z\in\C\mid |z|<1\}$. Let $g_X$ be the Poincar\'e metric, i.e., the complete K\"ahler metric whose Gaussian curvature equals $-1$. We denote by $h_X$ the corresponding Hermitian metric on $K_X\rightarrow X$ and by $\omega_X$ the K\"ahler form. Note that $h_X$ and $\omega_X$ satisfy the following:
\begin{align*}
\inum F_{h_X}=\omega_X.
\end{align*}
We choose the reference metric $h_\refe$ as $h_X$. First, following \cite[Section 5]{LM1}, we prove the existence of a complete solution to equation (\ref{phi}) under the assumption that $e^\varphi$ is smooth. We set a constant $M_\varphi$, which may be infinite, as follows:
\begin{align*}
M_\varphi\coloneqq \sup_X e^{\varphi}.
\end{align*}
For each $0<\epsilon<1$, let $\D_\epsilon\coloneqq \{z\in\C\mid |z|<1-\epsilon\}$ be the open disc of radius $1-\epsilon$ centered at $0$. Then we define a family of open subsets $(Y_i)_{i\geq 2}$ of $X=\D$ by $Y_i\coloneqq \D_{1/i}$ for $i=2,3,\dots$. For each $i=2,3,\dots$, we denote by $\partial Y_i$ the boundary of $Y_i$. Given each $i=2,3,\dots$ and each diagonal continuous Hermitian metric $\partial h^{(i)}=(\partial h_1^{(i)},\dots, \partial h_r^{(i)})$ on $\partial Y_i$, there uniquely exists a solution $h=(h_1,\dots, h_r)$ to equation (\ref{phi}) of class $C^{1,\alpha}$ for any $\alpha\in (0,1)$, such that $\lim_{z\to\zeta}h(z)=\partial h^{(i)}(\zeta)$ for any $\zeta\in \partial Y_i$ (cf. \cite{Miy3}). From $h^{(i)}$, we define Hermitian metrics $H_j^{(i)}$ on $K_X^{-1} \to X$ for each $j=1,\dots, r-1$ by
\begin{align*}
H_j^{(i)}\coloneqq h_j^{(i)-1}\otimes h_{j+1}^{(i)},
\end{align*}
and define $H_0^{(i)}=H_r^{(i)}$ by
\begin{align*}
H_0^{(i)} = H_r^{(i)} \coloneqq e^{r\varphi} h_\refe^{-r}|_{Y_i} \otimes h_r^{(i)-1} \otimes h_1^{(i)}.
\end{align*}

\begin{lemm}\label{cmphi}{\it Suppose that $e^\varphi$ is smooth and a bounded function. Then there exists a positive constant $C_{M_\varphi} $ depending only on $M_\varphi$ and $r$, such that for each $i = 1, 2, \dots$, by choosing the boundary condition appropriately, a solution $h^{(i)} = (h_1^{(i)}, \dots, h_r^{(i)})$ to the Dirichlet problem for equation (\ref{phi}) on $Y^{(i)}$ satisfies the following estimate:
\begin{align}
\frac{1}{2}\omega_X \leq \vol(H_j^{(i)}) \leq C_{M_\varphi} \omega_X \ \text{for each $j = 1, \dots, r-1$}. \label{cmphi}
\end{align}
}
\end{lemm}
\begin{proof} 
Let $C_1$ and $C_2$ be constants as in inequality (\ref{C1C2}). We set $C_{M_\varphi}$ as follows:
\begin{align*}
C_{M_\varphi}\coloneqq \max\left\{\frac{C_2M_\varphi+1}{C_1}, \frac{r-1}{2}+2^{r-1}M_\varphi^r\right\}. 
\end{align*}
For each $i = 1, 2, \dots$, we define a Hermitian metric $\partial h^{(i)} \coloneqq (\partial h^{(i)}_1, \dots, \partial h^{(i)}_r)$ on the boundary $\partial Y^{(i)}$ of $Y^{(i)}$ as the unique one such that $\partial h_j^{(i)\,-1} \otimes \partial h_{j+1}^{(i)} =\frac{1}{2}h_X^{-1}\left.\right|_{\partial Y^{(i)}}$ holds for all $j = 1, \dots, r - 1$. We show that the Hermitian metrics \( H_1^{(i)}, \dots, H_{r-1}^{(i)} \), constructed from the solution \( h^{(i)} = (h_1^{(i)}, \dots, h_r^{(i)}) \) to the Dirichlet problem for equation (\ref{phi}) with boundary condition \( \lim_{z \to \zeta} h(z) = \partial h^{(i)}(\zeta) \), satisfy the estimate~(\ref{cmphi}). We first prove the estimate $\frac{1}{2}\omega_X \leq \vol(H_j^{(i)})$ for each $j=1,\dots, r-1$. For each $j=1,\dots, r-1$, $H_j^{(i)}$ satisfies the following:
\begin{align*}
\inum\partial\bar{\partial}\log(H_j^{(i)}\otimes h_X)&=-\inum F_{H_j^{(i)}}-\inum F_{h_X} \\
&=2\vol(H_j^{(i)})-\vol(H_{j-1}^{(i)})-\vol(H_{j+1}^{(i)})-\omega_X \\
&\leq 2\vol(H_j^{(i)})-\omega_X.
\end{align*}
Then from the boundary condition and the strong maximum principle, we have $\frac{1}{2}\omega_X \leq \vol(H_j^{(i)})$ on $Y^{(i)}$. We next prove the estimate $\vol(H_j^{(i)}) \leq C_{M_\varphi}\omega_X$. From inequality (\ref{C1C2}) for $H=M_\varphi h_X^{-1}$ and the strong maximum principle, one of the following holds:
\begin{align}
&0\geq C_1\sum_{j=1}^rH_j^{(i)}\otimes H^{-1}-C_2-M_\varphi^{-1}, \\
&2^{r-1}M_\varphi^{-1}+2^{r-1}e^{r\varphi}\left.\right|_{\partial Y^{(i)}}M_\varphi^{-1}\geq \sum_{j=1}^rH_j^{(i)}\otimes H^{-1}.
\end{align}
In either case, we have $\vol(H_j^{(i)})\leq \sum_{j=1}^r\vol(H_j^{(i)})\leq C_{M\varphi}\omega_X$ for each $j=1,\dots, r-1$. This implies the claim.
\end{proof}
Then following \cite[Proposition 5.6]{LM1}, we show the following:
\begin{lemm}\label{sandb}{\it Suppose that $e^\varphi$ is smooth and bounded. Then there exists a smooth complete solution to equation (\ref{phi}).
}
\end{lemm}
\begin{proof} 
From Lemma \ref{cmphi} and the interior $L^p$-estimate (cf. \cite[Theorem 9.11]{GT1}), for each $i\geq 2$, $1<p<\infty$, and $l\geq 0$, there exist positive constants $C_1$ and $C_2$ that are independent of $l$ such that
\begin{align*}
|\log h^{(i+1+l)}|_{L^p_2(Y_i)}\leq C_1(|\log h^{(i+1+l)}|_{L^p(Y_{i+1})}+|F_{h^{(i+1+l)}}|_{L^p(Y_{i+1})})\leq C_2,
\end{align*}
where the $L^p_2$-norm and the $L^p$-norm are measured by the Poincar\'e metric. Then by the weak completeness of the reflexive Banach spaces, for each $i\geq 2$, we can find a subseqence of $(h^{(i+1+l)})_{l\geq 0}$ that converges to a Hermitian metric $h$ on $Y_i$ in the $L^p_2(Y_i)$-weak sense. By using the diagonal trick, we can find a subsequence of $(h^{(k)})_{k\geq 1}$ that converges to an $L^p_{2,loc}$-Hermitian metric $h=(h_1,\dots, h_r)$ on $X$ in the $L^p_{2,loc}$-weak sense. The Hermitian metric $h$ solves equation (\ref{phi}) in the weak sense. By using the standard elliptic regularity theorem, $h$ is smooth and solves equation (\ref{phi}) in the strong sense. From Lemma \ref{cmphi}, $h$ satisfies the following estimate:
\begin{align}
\frac{1}{2}\omega_X \leq \vol(H_j) \leq C_{M_\varphi} \omega_X \ \text{for each $j=1,\dots, r-1$}, \label{1/2omega}
\end{align}
where for each $j=1,\dots, r-1$, we have defined $H_j$ as $H_j\coloneqq h_j^{-1}\otimes h_{j+1}$. In particlular, $h$ is complete. 
\end{proof}
Following \cite[Proposition 5.7 and Proposition 5.8]{LM1}, we drop the assumption that $e^\varphi$ is a bounded function.
\begin{lemm}\label{unbounded}{\it Suppose that $e^\varphi$ is smooth. Then there exists a smooth complete solution $h$ to equation (\ref{phi}).
}
\end{lemm}
\begin{proof} 
For each $i\geq 2$, let $h_c^{(i)}=(h_{c,1}^{(i)},\dots, h_{c,r}^{(i)})$ be the unique complete solution to equation (\ref{phi}) on $Y_i$ associated with $e^{-\varphi}h_\refe|_{Y_i}$. Then, for each $k > i \geq 2$ and each $j = 1, \dots, n$, the quantity $h_{c,j}^{(k)-1} \otimes h_{c,j}^{(i)}(z)$ tends to zero as $z$ approaches the boundary point of $Y_i$, since $h_c^{(k)}$ is an incomplete Hermitian metric on $Y_i $, and by (\ref{1/2omega}), $h_c^{(i)}$ is dominated by the complete hyperbolic metric as follows:
\begin{align*}
&h^{(i)}_{c, n-k}\leq 2^kh_{Y_i}^kh_{c,n} \ \text{for each $k=1,\dots, n-1$}, \\
&(h_{c,n}^{(i)})^{2n+2-r}\leq 2h_{Y_i},
\end{align*} 
where we denote by $h_{Y_i}$ the Hermitian metric on $K_{Y_i}\rightarrow Y_i$ induced by the complete hyperbolic metric on $Y_i$. In particular, $h_{c,j}^{(k)-1} \otimes h_{c,j}^{(i)}$ is bounded above for each $j=1,\dots, n$. Then from Lemma \ref{cands}, we have $h_{c,j}^{(i)}<h_{c,j}^{(k)}$ for each $i<k$ and each $j=1,\dots, n$. Since $h_{c}^{(i)}$ is bounded above by $Ch_X$ with some positive constant $C$ independent of $i$, we see that $(h^{(k)})_{k\geq i}$ is uniformly bounded above and below on $Y_i$. Then from the interior $L^p$-estimate (\cite[Theorem 9.11]{GT1}), for each $i\geq 2$, $1<p<\infty$, and $l\geq 0$, there exist positive constants $C_1$ and $C_2$ that are independent of $l$ such that
\begin{align*}
|\log h_c^{(i+1+l)}|_{L^p_2(Y_i)}\leq C_1(|\log h_c^{(i+1+l)}|_{L^p(Y_{i+1})}+|F_{h_c^{(i+1+l)}}|_{L^p(Y_{i+1})})\leq C_2,
\end{align*}
where the $L^p$-norm and the $L^p_2$-norms are measured by the Poincar\'e metric on $X$. Then by the weak completeness of the reflexive Banach spaces, for each $i\geq 2$, we can find a subsequence of $(h_c^{(i+1+l)})_{l\geq 0}$ that converges to a Hermitian metric $h$ on $Y_i$ in the $L^p_2(Y_i)$-weak sense. By using the diagonal trick, we can find a subsequence of $(h_c^{(k)})_{k\geq 2}$ that converges to an $L^p_{2,loc}$-Hermitian metric $h=(h_1,\dots, h_r)$ on $X$ in the $L^p_{2,loc}$-weak sense. The Hermitian metric $h$ solves equation (\ref{phi}) in the weak sense. By using the standard elliptic regularity theorem, $h$ is smooth and solves equation (\ref{phi}) in the strong sense. Since we have the estimate $h_j\leq Ch_X^{\frac{r-(2j-1)}{2}}$ with some positive constant $C$ for each $j=1,\dots, n$, the metric $h$ is complete. This completes the proof. 
\end{proof}
From the proof of Lemma \ref{unbounded}, we directly obtain the following lemma:
\begin{lemm}\label{khn}{\it Suppose that $e^\varphi$ is smooth. Then the unique complete solution $h=(h_1,\dots, h_r)$ to equation (\ref{phi}) associated with $\varphi$ satisfying $\det(h)=1$ satisfies the following estimate:
\begin{align}
&h_{n-k}\leq 2^kh_X^kh_n \ \text{for each $k=1,\dots, n-1$}, \\
&h_n^{2n+2-r}\leq 2h_X.
\end{align}
}
\end{lemm}
Then we prove Theorems \ref{main theorem 2} and \ref{main theorem 3}:
\begin{proof}[Proof of Theorems \ref{main theorem 2} and \ref{main theorem 3}] Let $(\varphi_\epsilon)_{0<\epsilon<1}$ be a family of weight functions satisfying conditions in Theorem \ref{main theorem 3}. From Theorem \ref{main theorem 1} and Lemma \ref{unbounded}, there uniquely exists a smooth complete solution $h_\epsilon=(h_{\epsilon,1},\dots, h_{\epsilon,r})$ to equation (\ref{phi}). From the uniqueness of the complete solution, by applying the standard argument (cf. \cite[Corollary 3.27]{LM1}), $h_\epsilon$ is real for each $0<\epsilon<1$. As we remarked in Example \ref{subandsuper}, for each $\epsilon^\prime<\epsilon$, $h_{\epsilon^\prime}$ is a supersolution to equation (\ref{phi}) associtated with $\varphi_\epsilon$. Also, from Lemma \ref{khn}, since $h_{\epsilon^\prime}$ is incomplete on $\D_\epsilon$, for each $j=1,\dots, n$, $h_{\epsilon,j}\otimes h_{\epsilon^\prime,j}^{-1}(z)$ tends to zero as $z$ goes to the boundary of $\D_\epsilon$. In particular, for each $j=1,\dots, n$, $h_{\epsilon,j}\otimes h_{\epsilon^\prime,j}^{-1}$ is a bounded function on $\D_\epsilon$. Therefore from Lemma \ref{cands}, we obtain $h_{\epsilon,j}\otimes h_{\epsilon^\prime,j}^{-1}<1$ for each $j=1,\dots, n$. For each $i=2,3,\dots$, we denote by $h^{(i)}=(h^{(i)}_1,\dots, h^{(i)}_r)$ the complete solution $h_{1/i}$ defined on $Y_i=\D_{1/i}$. Then from the interior $L^p$-estimate (\cite[Theorem 9.11]{GT1}), for each $i\geq 2$, $1<p<\infty$, and $l\geq 0$, there exist positive constants $C_1$ and $C_2$ that are independent of $l$ such that
\begin{align*}
|\log h^{(i+1+l)}|_{L^p_2(Y_i)}\leq C_1(|\log h{(i+1+l)}|_{L^p(Y_{i+1})}+|F_{h^{(i+1+l)}}|_{L^p(Y_{i+1})})\leq C_2,
\end{align*}
where the $L^p$-norm and the $L^p_2$-norms are measured by the Poincar\'e metric on $X$. Then by the weak completeness of the reflexive Banach spaces, for each $i\geq 2$, we can find a subsequence of $(h^{(i+1+l)})_{l\geq 0}$ that converges to a Hermitian metric $h$ on $Y_i$ in the $L^p_2(Y_i)$-weak sense. By using the diagonal trick, we can find a subsequence of $(h^{(k)})_{k\geq 2}$ that converges to an $L^p_{2,loc}$-Hermitian metric $h=(h_1,\dots, h_r)$ on $X$ in the $L^p_{2,loc}$-weak sense. By definition, $h$ solves equation (\ref{phi}) associated with $\varphi$ in the weak sense. Since $h^{(i)}$ is real for each $i$, $h$ is also real. By the sobolev embedding theorem and the elliptic regularity theorem, for each $j=1,\dots, n$, $h_j$ is of class $C^{2j-1,\alpha}$ for any $\alpha\in(0,1)$. Since we have the estimate $h_j\leq Ch_X^{\frac{r-(2j-1)}{2}}$ for each $j=1,\dots, n$, the metric $h$ is complete. This proves Theorem \ref{main theorem 2}. Since we have $h_{\epsilon,j}\otimes h_{\epsilon^\prime,j}^{-1}<1$ for each $0<\epsilon^\prime<\epsilon<1$ and each $j=1,\dots, n$, $(h_\epsilon)_{0<\epsilon<1}$ monotonically converges to $h$ as $\epsilon\searrow 0$. This proves Theorem \ref{main theorem 3}. 

\end{proof}




\begin{thebibliography}{99}
\bibitem{AF1}E. Aldrovandi and G. Falqui, {\it Geometry of Higgs and Toda fields on Riemann surfaces}, J. Geom. Phys. 17 (1995), no. 1, 25-48.
\bibitem{Bar1} D. Baraglia, {\it $G_2$ geometry and integrable systems}, arXiv:1002.1767 (2010).
\bibitem{BCHM1} T. Bayraktar, D. Coman, H. Herrmann, and G. Marinescu, {\it A survey on zeros of random holomorphic sections}, arXiv:1807.05390 (2018).
\bibitem{BB1} R. Berman and S. Boucksom, {\it Growth of balls of holomorphic sections and energy at equilibrium}, Inventiones mathematicae, 181(2), 337-394 (2010).
\bibitem{CY1} S. Y. Cheng and S. T. Yau, {\it Differential equations on Riemannian manifolds and their geometric applications}, Communications on Pure and Applied Mathematics 28.3 (1975): 333-354.
\bibitem{CT1} H. I. Choi and A. Treibergs, {\it Gauss maps of spacelike constant mean curvature hypersurfaces of Minkowski space}, Journal of differential geometry 32.3 (1990): 775-817.
\bibitem{Col1} B. Collier, {\it Finite order automorphisms of Higgs bundles: theory and application}, (Doctoral dissertation, University of Illinois at Urbana-Champaign).
\bibitem{DL1} S. Dai, and Q. Li, {\it Minimal surfaces for Hitchin representations}, Journal of Differential Geometry 112.1 (2019): 47-77.
\bibitem{DL2} S. Dai and Q. Li, {\it On cyclic Higgs bundles}, Math. Ann. 376 (2020), no. 3-4, 1225-1260.
\bibitem{DL3} S. Dai and Q. Li, {\it Bounded differentials on the unit disk and the associated geometry}, Transactions of the American Mathematical Society (2024).
\bibitem{Don1} S. K. Donaldson, {\it Boundary value problems for Yang-Mills fields}, Journal of geometry and physics 8.1-4 (1992): 89-122.
\bibitem{DW1} D. Dumas and M. Wolf, {\it Polynomial cubic differentials and convex polygons in the projective plane}, Geometric and Functional Analysis 25.6 (2015): 1734-1798.
\bibitem{FKMR1} O. El-Fallah, K. Kellay, J. Mashreghi, and T. Ransford, {\it A primer on the Dirichlet space} (Vol. 203). Cambridge University Press (2014).
\bibitem{Gal1} G. Gallavotti, {\it Statistical mechanics: A short treatise}, Springer Science \& Business Media, 1999.
\bibitem{GT1} D. Gilbarg and N. Trudinger, {\it Elliptic partial differential equations of second order}, Vol. 224. No. 2. Berlin: springer, 1977.
\bibitem{GZ1} V. Guedj and A. Zeriahi, {\it Degenerate Complex Monge-Amp\`ere Equations}, EMS Tracts Math. 26, European Mathematical Society (EMS), Z\"urich 2017.
\bibitem{GH1} M. A. Guest and N. K. Ho, {\it A Lie-theoretic description of the solution space of the $tt^\ast$-Toda equations}, Mathematical Physics, Analysis and Geometry 20.4 (2017): 24.
\bibitem{GL1} M. A. Guest and C.-S. Lin, {\it Nonlinear PDE aspects of the $tt^\ast$ equations of Cecotti and Vafa}, J. Reine Angew. Math. 689 (2014), 1-32.
\bibitem{Hit1} N. J. Hitchin, {\it The self-duality equations on a Riemann surface}, Proceedings of the London Mathematical Society 3.1 (1987): 59-126.
\bibitem{Hit2} N. J. Hitchin, {\it Lie groups and Teichm\"uller space}, Topology, 31(3), 449-473 (1992).
\bibitem{LT1} F. Labourie and J. Toulisse, {\it Quasicircles and quasiperiodic surfaces in pseudo-hyperbolic spaces}, Inventiones mathematicae 233.1 (2023): 81-168.
\bibitem{LTW1} F. Labourie, J. Toulisse, and M. Wolf, {\it Plateau problems for maximal surfaces in pseudo-hyperbolic spaces}, arXiv:2006.12190 (2020).
\bibitem{Li0} Q. Li, {\it Harmonic maps for Hitchin representations}, Geometric and Functional Analysis 29 (2019): 539-560.
\bibitem{Li1} Q. Li, {\it On the uniqueness of vortex equations and its geometric applications}, The Journal of Geometric Analysis 29 (2019): 105-120.
\bibitem{LM1}Q. Li and T. Mochizuki, {\it Complete solutions of Toda equations and cyclic Higgs bundles over non-compact surfaces}, International Mathematics Research Notices 2025.7 (2025): rnaf081..
\bibitem{LM2} Q. Li and T. Mochizuki, {\it Isolated singularities of Toda equations and cyclic Higgs bundles}, arXiv:2010.06129 (2020).
\bibitem{Lof1} J. C. Loftin, {\it Flat metrics, cubic differentials and limits of projective holonomies}, Geometriae Dedicata, 128 (2007), 97-106.
\bibitem{Miy2} N. Miyatake, {\it Generalizations of Hermitian-Einstein equation of cyclic Higgs bundles, their heat equation, and inequality estimates}, arXiv:2301.01584.
\bibitem{Miy3} N. Miyatake, {\it Cyclic Higgs bundles, subharmonic functions, and the Dirichlet problem}, Annals of Global Analysis and Geometry 67.1 (2025): 7.
\bibitem{Miy4} N. Miyatake, {\it Shannon entropy for harmonic metrics on cyclic Higgs bundles}, arXiv:2410.08571.
\bibitem{Miy5} N. Miyatake, {\it Shannon entropy for harmonic metrics on cyclic Higgs bundles II}, preprint.
\bibitem{Moc0} T. Mochizuki, {\it Harmonic bundles and Toda lattices with opposite sign I}, arXiv:1301.1718.
\bibitem{Moc1} T. Mochizuki, {\it Harmonic bundles and Toda lattices with opposite sign II.}, Communications in Mathematical Physics 328 (2014): 1159-1198.
\bibitem{Nie1} X. Nie, {\it Poles of cubic differentials and ends of convex $\mathbb {RP}^ 2$-surfaces}, Journal of Differential Geometry 123.1 (2023): 67-140.
\bibitem{Omo1} H. Omori, {\it Isometric immersions of Riemannian manifolds}, Journal of the Mathematical Society of Japan 19.2 (1967): 205-214.
\bibitem{Ran1} T. Ransford, {\it Potential theory in the complex plane}, No. 28. Cambridge university press, 1995.
\bibitem{ST1} E. B. Saff and V. Totik, {\it Logarithmic potentials with external fields}, Vol. 316. Springer Science \& Business Media, 2013.
\bibitem{Sh1} C. E. Shannon, {\it A mathematical theory of communication}, The Bell system technical journal, 27(3), 379-423 (1948).
\bibitem{SZ1} B. Shiffman and S. Zelditch, {\it Stochastic K\" ahler geometry: from random zeros to random metrics}, arXiv:2303.11559 (2023).
\bibitem{Sim1} C. T. Simpson, {\it Constructing variations of Hodge structure using Yang-Mills theory and applications to uniformization}, Journal of the American Mathematical Society (1988): 867-918.
\bibitem{Sim2} C. T. Simpson, {\it Harmonic bundles on noncompact curves}, Journal of the American Mathematical Society 3.3 (1990): 713-770.
\bibitem{Sim3} C. T. Simpson, {\it Katz's middle convolution algorithm}, arXiv preprint math/0610526. (2006).
\bibitem{Yau1} S. T. Yau, {\it Harmonic functions on complete Riemannian manifolds}, Communications on Pure and Applied Mathematics 28.2 (1975): 201-228.
\bibitem{Wan1} T. Y.-H. Wan, {\it Constant mean curvature surface, harmonic maps, and universal Teichm\"uller space}, J. Differential Geom.35 (1992), no.3, 643-657.
\bibitem{WA1} T. Y.-H. Wan and T. K.-K. Au, {\it Parabolic constant mean curvature spacelike surfaces}, Proceedings of the American Mathematical Society 120.2 (1994): 559-564.
\bibitem{Wol1} M. Wolf, {\it The Teichm\"uller theory of harmonic maps}, Journal of differential geometry 29.2 (1989): 449-479.
\end{thebibliography}

\noindent
E-mail address 1: natsuo.miyatake.e8@tohoku.ac.jp

\noindent
E-mail address 2: natsuo.m.math@gmail.com \\

\noindent
Mathematical Science Center for Co-creative Society, Tohoku University, 468-1 Aramaki Azaaoba, Aoba-ku, Sendai 980-0845, Japan.


\end{document}












